% ============================================================
% Viés Geográfico e Linguístico em LLMs para Política de Poluição do Ar
% Versão em português — acompanha latex/main.tex (inglês)
% ============================================================
% ESTA VERSÃO É MODULAR, como a inglesa. A anterior era um arquivo único de
% 1.615 linhas, e foi exatamente por isso que ficou dois meses defasada sem que
% ninguém percebesse: manter dois monólitos em paralelo é garantir que divirjam.
% Cada seção aqui espelha o arquivo correspondente em sections/, e os NÚMEROS são
% os mesmos, vindos do mesmo congelamento (data/processed/freeze_all_effects.json).
% Decimais com ponto, como na versão anterior, para não introduzir erro de
% conversão em algumas centenas de valores.
% ============================================================

\documentclass[11pt,a4paper]{article}

\usepackage[utf8]{inputenc}
\usepackage[T1]{fontenc}
\usepackage[brazil]{babel}

\usepackage[margin=2.5cm]{geometry}
\usepackage{setspace}
\onehalfspacing
\usepackage{lineno}
\linenumbers

\usepackage{amsmath,amssymb}
\usepackage{siunitx}
\sisetup{group-separator={.}, group-minimum-digits=4}

\usepackage{graphicx}
\graphicspath{{../figures/}}
\usepackage{xcolor}
\usepackage{booktabs}
\usepackage{longtable}
\usepackage{array}
\usepackage{caption}
\captionsetup{font=small,labelfont=bf}

\usepackage[authoryear,round]{natbib}
\bibliographystyle{apalike}

% URLs longas e nomes compostos estouram a margem; ver main.tex
\usepackage{url}
\def\UrlBreaks{\do\/\do\-\do\.\do\_}
\Urlmuskip=0mu plus 1mu
\hyphenation{Deep-Seek An-thro-pic Son-net Ge-mi-ni Kripp-en-dorff
             mul-ti-lin-gue pre-es-pe-ci-fi-ca-do re-tro-tra-du-cao
             ad-mi-nis-tra-cao}
\setlength{\emergencystretch}{3em}

\usepackage[colorlinks=true,
            linkcolor=blue!50!black,
            citecolor=blue!50!black,
            urlcolor=blue!50!black]{hyperref}

\newcommand{\placeholder}[1]{\textcolor{red!80!black}{\textbf{[#1]}}}

% ============================================================

\title{De quem é a regulação que o modelo conhece?\\
Confiabilidade da informação regulatória recuperada por LLMs para a gestão
ambiental pública em 25 países}

\author{
  Lucas Rover$^{1,*}$ \\
  Vitor de Melo Dominski$^{2}$ \\
  Anibal Tavares de Azevedo$^{3}$ \\
  Eduardo Tadeu Bacalhau$^{4}$ \\
  Yara de Souza Tadano$^{1}$ \\
  \\
  \small $^1$ Universidade Tecnológica Federal do Paraná (UTFPR), Brasil\\
  \small $^2$ Faculdade Descomplica, Brasil\\
  \small $^3$ Universidade Estadual de Campinas (Unicamp), Brasil\\
  \small $^4$ Universidade Federal do Paraná (UFPR), Brasil\\
  \small $^*$ Autor correspondente: \texttt{lucasrover@alunos.utfpr.edu.br}\\
  \\
  \small ORCID: L.R. 0000-0001-6641-9224 $\cdot$ V.M.D. 0009-0001-0209-9089
  $\cdot$ A.T.A. 0000-0003-1678-7795 $\cdot$ E.T.B. 0000-0002-3936-0375
  $\cdot$ Y.S.T. 0000-0002-3975-3419
}

\date{\today}

\begin{document}

\maketitle

% Bloco unico, como a Elsevier exige, e ate 250 palavras. Cada "fix" declarado
% com o status que a inferencia no nivel do pais sustenta (parecer de painel,
% rodada 2, P0).
\begin{abstract}

Um gestor ambiental municipal pergunta a um modelo de linguagem qual é o padrão
vinculante de PM\textsubscript{2,5} em seu país, e nada na resposta diz se ela
corresponde ao registro oficial. Propomos um protocolo para essa conferência e
relatamos o que ele encontra. As respostas são pontuadas contra registros
oficiais, por código sempre que a resposta correta é um valor publicado e por um
painel de juízes de três fornecedores nos demais casos, em 25 países, 14
configurações de acesso e quatro idiomas (\ncellspt{} células pontuadas).
Nenhuma das três mitigações que um órgão consegue implantar sem construir nada
funciona, e a mais intuitiva prejudica. Perguntar no idioma do próprio país
reduz a acurácia em \nnativa{} pontos percentuais (modelo misto com intercepto
aleatório por país, $p=\nhtwomixedcountryp$; negativa em \nhtwocneg{} de nove
países) e torna o modelo \nverifdiffabs{} pontos menos propenso a devolver
qualquer valor conferível contra o registro. Declarar o papel oficial equivale a
não declarar nada, em até dois pontos percentuais. O ajuste regional de
instrução em 7B não melhora o modelo-base de que partiu. A acurácia é mais baixa
na tarefa com consequência administrativa, recordar o valor vinculante
(\ntasktoneshort{} contra \nfloorrestshort{} para síntese e recomendação), e um
terço das respostas erradas cita um valor da própria escada oficial do país. Os
países do Sul Global devolvem o valor do registro \ncodegapabs{} pontos menos
vezes que os do Norte com o país como unidade; as estimativas pontuais são
maiores onde a resposta é um fato nacional, mas todo intervalo por tarefa
inclui~1.

\end{abstract}


% Cinco, e nenhuma repete o título — ver a nota em main.tex.
\noindent\emph{Palavras-chave:} auditoria de algoritmos; viés geográfico; viés
multilíngue; governo digital; Sul Global

\vspace{1em}

\section{Introdução}
\label{sec:intro}

Os grandes modelos de linguagem tornaram-se infraestrutura de trabalho da
administração pública em poucos anos. Um levantamento da OCDE sobre funções
centrais de governo documenta a adoção de IA generativa em formulação de
políticas e prestação de serviços, incluindo assistentes que ajudam assessores
do serviço público a recuperar informação e as fontes por trás
dela~\citep{oecd2024governing}. A política de poluição do ar é um domínio em que
isso tem consequência direta para a saúde pública. O material particulado fino
está entre os principais fatores de risco ambiental para mortalidade prematura
no mundo~\citep{gbd2019risk,who2021aqg}, e a carga recai de forma
desproporcional sobre o Sul Global. Quando um gestor ambiental pergunta a um
modelo sobre um padrão de qualidade do ar, um prazo de atendimento ou a
evidência por trás de uma intervenção, uma resposta imprecisa pode se propagar
para um parecer regulatório, uma decisão de compra ou uma resposta de emergência
durante um episódio de poluição.

Este artigo pergunta se os modelos respondem a essas questões com a mesma
confiabilidade para países do Sul Global e do Norte Global, e se a forma de
enquadrar a consulta muda a resposta. As duas premissas têm apoio. O viés
factual geográfico é mensurável: as taxas de erro são $1,5\times$ maiores para
países da África Subsaariana do que para a América do Norte em 20
modelos~\citep{moayeri2024worldbench}, as avaliações dos modelos sobre temas
subjetivos acompanham condições socioeconômicas até
$\rho=0,70$~\citep{manvi2024llm}, e os modelos privilegiam sistematicamente
afirmações sobre o Norte Global~\citep{mirza2024global}. A assimetria linguística
por trás disso é estrutural: os idiomas da maioria dos países do Sul Global
ocupam as classes mal servidas da hierarquia de recursos~\citep{joshi2020state},
e grandes corpora derivados da Internet ``super-representam pontos de vista
hegemônicos e codificam vieses potencialmente danosos a populações
marginalizadas''~\citep{bender2021stochastic}. A produção decolonial lê essa
assimetria como uma infraestrutura epistêmica desigual reproduzida pelos
sistemas de IA~\citep{mohamed2020decolonial}.

Três lacunas organizam o estudo. Nenhum benchmark visa informação regulatória
com gabarito oficial: as avaliações existentes usam recordação de indicadores do
Banco Mundial, avaliações subjetivas, checagem jornalística ou conhecimento
cultural cotidiano, e nenhuma mede as tarefas que um gestor delega, como
enunciar um padrão vinculante, recuperar uma concentração medida oficialmente,
resumir a evidência de carga de doença, identificar instrumentos legais e órgãos
executores ou recomendar uma resposta operacional. O enquadramento por persona
não foi examinado como alavanca experimental, embora os praticantes prefixem um
papel rotineiramente e os fornecedores o recomendem (Seção~\ref{sec:related});
se ele reduz a lacuna, ao direcionar para respostas conformes à regulação, ou a
amplia, ao convidar à fabricação confiante, não foi testado em desenho
pré-especificado. E o mecanismo de representação no corpus costuma ser afirmado
em vez de medido contra confundidores, e raramente é separado em seus dois
canais, o tamanho do corpus do idioma~\citep{xue2021mt5,abadji2022oscar} e a
amplitude da cobertura do país, distinção que importa aqui porque a lacuna
geográfica é medida em inglês.

\subsection*{Objetivos e escopo}

O estudo tem três objetivos delimitados. O primeiro é estabelecer um protocolo
de auditoria de informação regulatória. Seu objeto é a concordância entre a
resposta de um modelo e o valor publicado no registro oficial, para tarefas que
um órgão ambiental delega; ele não mede capacidade geral do modelo, qualidade de
raciocínio ou utilidade para outros fins. O protocolo é a contribuição que se
transfere, e o domínio é onde o demonstramos. O segundo é testar as mitigações
disponíveis a uma administração. Um órgão público diante de respostas ruins
sobre o próprio país pode declarar o papel oficial no prompt, perguntar no
idioma nacional ou rodar um modelo aberto com ajuste regional, e testamos as
três como fatores de desenho. O modelo regional que pudemos rodar é um ajuste
comunitário de uma base de 7B em português brasileiro, do tipo que um órgão
consegue baixar; o achado diz respeito a essa classe, e o protocolo liberado
permite testar qualquer outro candidato. Não testamos recuperação aumentada,
ajuste fino ou verificação com humano no circuito, cada um dos quais exige
capacidade de engenharia que as correções no nível do prompt não exigem. Todas as consultas
são respondidas a partir da memória paramétrica do modelo, sem busca na web nem
recuperação de documentos, e por isso o título fala em informação recordada, e
não recuperada. O terceiro objetivo é localizar onde a confiabilidade difere e,
até onde o desenho permite, por quê; testamos a explicação por representação no
corpus e não reivindicamos identificação causal.

Ao longo do texto, a unidade de análise é o modelo tal como servido, isto é,
pesos, fornecedor e infraestrutura de serviço alcançados por um endpoint fixado,
porque o gestor público interage com essa configuração e qualquer lacuna medida é
propriedade dela. Hipóteses, amostragem e plano de análise foram fixados antes
da coleta~\citep{nosek2018preregistration} para 15 países e depois estendidos a
25; o plano nunca foi depositado publicamente, portanto não reivindicamos
pré-registro, e o estudo é relatado como exploratório do início ao fim.

\subsection*{Contribuições deste trabalho}

\begin{enumerate}\setlength{\itemsep}{2pt}
\item \emph{Um protocolo de auditoria para informação regulatória.} A pergunta
  que uma administração enfrenta é se o valor que o modelo devolve corresponde
  ao registro oficial. Construímos cinco tipos de tarefa a partir do trabalho
  cotidiano de um órgão ambiental (Tabela~\ref{tab:tarefas}), cada um atrelado a
  uma fonte oficial, de modo que a correção é adjudicável em vez de uma questão
  de impressão de especialista, e o protocolo se transfere a qualquer domínio
  regulado em que a resposta correta é publicada.
\item \emph{As três correções de localização a que os órgãos recorrem,
  testadas em vez de presumidas.} Enquadramento por persona, prompt em idioma
  nativo e modelo com ajuste regional são, cada um, uma decisão de compra ou de
  fluxo de trabalho. Nenhum deles funciona, e falham de maneiras diferentes:
  perguntar no idioma do próprio país reduz a acurácia em $\nnativa$ pontos
  percentuais, de forma mais acentuada para o idioma de corpus mais escasso; o
  enquadramento de papel equivale a nenhum enquadramento em até dois pontos
  percentuais; e o modelo com ajuste regional recorda o valor vinculante com a
  mesma frequência do modelo-base de que partiu.
\item \emph{Um desenho que separa quem é mal servido de por quê, e relata onde
  não consegue.} Os países são estratificados pela divisão desenvolvido/em
  desenvolvimento da UNCTAD, pela taxonomia de recursos linguísticos de Joshi e
  pelo desenvolvimento, de modo que explicações geográficas, linguísticas e
  econômicas possam ser distinguidas. Entre países, cobertura e desenvolvimento
  são colineares; dentro dos países, o que prevê o déficit de especificidade no
  nível do país é o desenvolvimento, e não a cobertura enciclopédica que o
  protocolo pré-especificou, e os dois não se separam nesta amostra. Documentamos também
  que o proxy de corpus pré-especificado atribui a Wikipédia inglesa a todo
  país da amostra com inglês como língua oficial, de modo que ele acompanha a
  história colonial e não o tamanho do corpus.
\item \emph{Gabarito adjudicado por código, e o que isso muda.} Reconstruir
  duas chaves de tarefa a partir de registros oficiais e mover a comparação de
  faixa para código reduziu o gradiente de desenvolvimento de $\rho=0,51$ para
  $\rho=\nrho$ e mais que dobrou a penalidade do idioma nativo nas mesmas
  respostas; mover a terceira tarefa com valor de registro de um juiz para
  código removeu um confundimento entre juiz e tier. Nesta literatura, a chave
  de gabarito, e não o tamanho da amostra, é a restrição que vincula.
\item \emph{Recursos abertos.} Protocolo, prompts, registros de gabarito com
  suas fontes oficiais, código de análise e o registro de desvios são liberados
  sob CC-BY-4.0 e MIT, com pipeline de reprodução em um comando.
\end{enumerate}

\section{Trabalhos Relacionados}
\label{sec:related}

A literatura prévia corre por quatro fios: a medição do viés factual
geográfico, extensões multilíngues e culturais, tentativas de mitigação, e o uso
de LLMs em contextos de decisão de política e saúde.

\subsection{Viés factual geográfico e extensões multilíngues}

\citet{manvi2024llm} estabeleceram que LLMs de fronteira carregam vieses
sistemáticos contra regiões de condições socioeconômicas mais baixas, com
correlações de até 0,70 entre avaliações geradas pelo modelo e um proxy dessas
condições; o espaço de desfecho são avaliações subjetivas, o que deixa em aberto
se o mesmo viés aparece na recordação objetiva de fatos regulatórios.
\citet{moayeri2024worldbench} quantificaram a recordação factual contra
indicadores do Banco Mundial em 20 LLMs, documentando erro relativo absoluto
$1{,}5\times$ maior para a África Subsaariana do que para a América do Norte; o
desenho é limitado pelo espaço de indicadores macroeconômicos, enquanto um
gestor ambiental faz perguntas mais estreitas e específicas de fonte, cujo
gabarito vive em conselhos ambientais e agências de monitoramento.
\citet{mirza2024global} mostraram que versões mais novas do GPT-4 não melhoram
monotonicamente e que o privilégio de afirmações sobre o Norte Global persiste.

A literatura de sondagem multilíngue estabeleceu cedo que o que um modelo
recorda depende do idioma da consulta. \citet{jiang2020xfactr} sondaram fatos
em formato cloze em 23 idiomas e encontraram acurácia abaixo de 15\% mesmo em
inglês e espanhol e abaixo de 5\% em marathi e iorubá, com metade dos fatos
corretamente recordados recordados em um único idioma; \citet{kassner2021mlama}
traduziram os mesmos benchmarks para 53 idiomas e encontraram o mBERT abaixo
de 60\% do seu desempenho em inglês em 32 deles, um viés para entidades do
idioma da consulta e nenhuma relação clara entre o tamanho da Wikipédia de um
idioma e sua recordação. Esses resultados dizem respeito a codificadores
mascarados de 2020; a pergunta que levantam, se a recordação em um idioma
acompanha o corpus desse idioma, é a que H2 e H4 fazem a modelos generativos.
\citet{myung2024blend} cobriram 16 países e 13 idiomas em conhecimento
cultural cotidiano, constatando que os LLMs têm melhor desempenho
nos idiomas nativos de culturas de alto recurso e pior nos de baixo recurso, o
achado que motiva H2 e se conecta à hierarquia de recursos de
\citet{joshi2020state}. Benchmarks regionais proliferaram:
BRoverbs~\citep{broverbs2025}, TiEBe~\citep{almeida2025tiebe},
ALBA~\citep{alba2026} e AMALIA~\citep{amalia2026} documentam falhas concretas
em tarefas culturalmente situadas, mas não tratam do uso aplicado em política
pública. O trabalho de construção de corpora em português é explícito em que o
esforço se concentrou no inglês: \citet{almeida2025portuguese} montam um corpus
de 120 bilhões de tokens em português porque pipelines específicos por idioma
são necessários para alcançar qualidade industrial. Essa assimetria na
quantidade de texto por idioma é o que H4 testa.

Esses estudos estabelecem que o fenômeno existe e é mensurável. Nenhum isola um
único domínio aplicado com gabarito oficial verificável, e nenhum trata o
enquadramento do prompt como fator manipulável.

\subsection{Mitigação, persona e condicionamento por papel}

\citet{opuszko2026unraveling} constataram que o raciocínio ajuda no agregado,
mas não elimina as disparidades regionais, e \citet{kerche2026silicon}
propuseram uma tipologia baseada em lugar que situa o viés geográfico num
panorama sociotécnico mais amplo. Ambos são conceituais e agregados.

O condicionamento por papel é o que os fornecedores instruem os próprios
usuários a fazer. A Anthropic documenta a definição de papel no prompt de
sistema~\citep{anthropic2026systemprompts}, o guia da OpenAI abre a estrutura
de mensagem recomendada com um bloco de identidade~\citep{openai2026promptengineering},
a documentação do Gemini, da Google, pede que definições de papel, que ela
chama de persona, sejam colocadas na instrução de
sistema~\citep{google2026promptingstrategies}, o exemplo padrão da xAI inicia
toda conversa com um prompt de sistema que atribui um
papel~\citep{xai2026chatguide}, e o padrão está catalogado na literatura de
prompting~\citep{white2023prompt}. A evidência sistemática sobre se isso ajuda
é negativa: em quatro famílias de LLM e 2.410 perguntas factuais, personas de
especialista não melhoraram a acurácia em relação a um controle sem persona e
a reduziram ligeiramente em 162 papéis~\citep{zheng2024helpful}. Um estudo
independente com seis modelos em questões de múltipla escolha de nível de
pós-graduação chega à mesma conclusão: casar uma persona de especialista com o
tipo de problema não teve efeito significativo no
desempenho~\citep{mollick2025personas}. Se o condicionamento por persona ajuda
ou prejudica a confiabilidade factual em domínios de alto risco é, portanto,
questão em aberto e, até onde sabemos, não testada para equidade geográfica.
Tratamos a persona como fator experimental pré-especificado e testamos tanto a
hipótese de conformidade à norma (a persona reduz a lacuna) quanto a hipótese
do vácuo de persona (a persona a amplia).

\subsection{LLMs em política, saúde e na administração}

\citet{lecoz2025policymaking} avaliaram recomendações de política de LLMs
para a situação de rua em quatro cidades (Barcelona, Joanesburgo, South Bend e
Macau) e leram o padrão como ``uma rigidez cega ao contexto e um viés potencial
baseado no discurso online publicado sobre a situação de rua nas regiões do
mundo com dados super-representados nos LLMs'', em contraste com especialistas
que ajustaram as decisões a realidades codificadas localmente. Trata-se de um
domínio e um tipo de tarefa, que generalizamos para cinco tipos de tarefa
cobrindo recordação, síntese e recomendação. O precedente mais próximo num
domínio regulado é \citet{dahl2024legal}, que fizeram a quatro modelos públicos
perguntas verificáveis sobre casos federais dos Estados Unidos e encontraram
taxas de alucinação entre 58\% e 88\%, mais altas para as cortes inferiores e
as jurisdições menos proeminentes, e mais altas de novo para os casos mais
antigos e os mais recentes, de modo que o conhecimento dos modelos fica atrás
da doutrina em vigor; nosso piso de recordação e a análise da escada na
Seção~\ref{sec:res:t1} são a contraparte ambiental desses dois gradientes. Na
saúde, \citet{kozlakidis2026medical} argumentam que alucinações carregam risco agudo
em ambientes regulatórios com capacidade técnica limitada, onde o descompasso
de idioma e cultura aumenta o erro, um risco que encontra a carga de
mortalidade documentada pelas estimativas do Global Burden of
Disease~\citep{gbd2019risk}.

A questão da implantação já não é hipotética. \citet{kim2026genai} estima
efeitos de produtividade burocrática da IA generativa de forma
quase-experimental, e \citet{robinson2026opensource} documenta que escolher
qual modelo um órgão roda é hoje uma decisão explícita de compra entre opções
de pesos abertos e hospedadas pelo fornecedor. Nossos achados no nível do
modelo falam a essa decisão: os modelos de pesos abertos, como grupo, ficam
$\nhfive$~pontos abaixo dos fechados no composto, ainda que um modelo de pesos
abertos (DeepSeek-V3) lidere a lista, e o modelo português com ajuste regional
seja o mais fraco dos catorze. Dois outros fios enquadram por que uma resposta
errada importa administrativamente. \citet{choroszewicz2026crumple} mostra que
ferramentas de apoio com IA generativa posicionam os profissionais da linha de
frente como ``zonas de deformação moral'', absorvendo a responsabilidade por
falhas originadas a montante, e \citet{cordella2024regulating}, ao analisar a
primeira medida vinculante que um regulador impôs a um serviço de IA generativa
(a intervenção italiana sobre o ChatGPT), argumentam que arcabouços
tecnologicamente neutros ignoram o que é específico desses sistemas e que os
reguladores precisam engajar diretamente seu funcionamento técnico. Sobre a
recepção, \citet{aoki2024explainable} testam se o tipo de explicação muda a
acurácia, a justiça e a confiabilidade percebidas entre os afetados.

Nada disso estabelece se a resposta subjacente está correta, nem se a correção
se distribui de forma uniforme entre os Estados que dependem dela.
Confiabilidade percebida e acurácia medida são grandezas diferentes, e um
sistema pode ser bem explicado, bem governado, bem recebido e errado. Essa é a
lacuna que este estudo enfrenta, num domínio em que a resposta correta está
publicada num registro oficial e pode ser conferida.

\subsection{Posicionamento do presente trabalho}

O desenho difere do trabalho prévio em quatro aspectos: um único domínio
aplicado de alto risco com gabarito oficial verificável, em vez de trivialidades
numéricas ou conhecimento cultural; persona como fator intra-sujeito
pré-especificado, em vez de escolha de enquadramento não controlada; amostragem
de países orientada por teoria em eixos independentes (a divisão
desenvolvido/em desenvolvimento da UNCTAD, recursos linguísticos de Joshi,
desenvolvimento), em vez de conveniência regional; e um teste explícito do
mecanismo de representação no corpus com análise de sensibilidade. Nenhum
estudo anterior combina um domínio regulado de saúde pública, uma manipulação
de persona e um teste de mecanismo pré-especificado.

\section{Métodos}
\label{sec:metodos}

\subsection{Desenho da pesquisa}

Especificamos um experimento de benchmark pré-especificado, fatorial e multipaís,
quantificando viés geográfico e modulado por persona em 14 modelos de linguagem de
grande porte em tarefas de política de poluição do ar. O desenho pré-especificado
cruza 15 países (estratificados em três eixos teoricamente fundamentados), 14 LLMs
abrangendo cinco camadas de acesso, um domínio, 5 tipos de tarefa e 2 condições de
persona, com 2 replicações estocásticas por célula; uma extensão post hoc
(Seção~\ref{sec:expansao}) amplia a amostra para 25 países. O conjunto de estímulos
cruza as cinco tarefas e as duas personas de cada país em inglês, mais uma versão
em idioma nativo dentro do país para os nove países com idioma nativo alvo,
pontuado no composto $[0,1]$. O protocolo fixa, antes da coleta, as hipóteses, os
modelos primários, o protocolo de persona, os critérios de exclusão e as correções
para múltiplas comparações. Uma coleta apenas do Brasil, executada sob um desenho
anterior de três domínios, é reclassificada como piloto estendido e relatada
separadamente; ela informou a calibração dos prompts, mas não contribui com
observações confirmatórias.

\subsection{Seleção de países}

Os países foram selecionados por amostragem teórica~\citep{patton2015qualitative}
triangulada em três classificações independentes: a distinção Norte/Sul Global,
operacionalizada pela classificação de economias desenvolvidas/em desenvolvimento
da UNCTAD~\citep{unctad2024classification}, que carrega o arcabouço da
colonialidade do saber que motiva o estudo~\citep{mohamed2020decolonial}; a
taxonomia de seis classes de representação de recursos linguísticos de Joshi et
al.~\citep{joshi2020state}; e os grupos de renda do Banco
Mundial~\citep{worldbank2025classification}.

A amostra pré-especificada de 15 --- Brasil, México, Argentina, Peru; Nigéria,
África do Sul, Quênia, Egito; Índia, Indonésia, Bangladesh, Filipinas; e Estados
Unidos, Alemanha, Japão como referências do Norte Global --- cobre quatro regiões
do mundo e quatro grupos de renda. O eixo de recursos linguísticos revelou-se muito
menos diferenciado do que se pretendia: os idiomas oficiais dominantes caem em
apenas três classes de Joshi, doze dos quinze compartilhando a Classe~5, e estender
para 25 não ajuda, já que dezoito são Classe~5. Isso é propriedade do mundo, e não
do nosso sorteio, porque os idiomas oficiais da maioria dos Estados são eles
próprios de alto recurso; mas significa que o eixo linguístico tem variância
pequena demais para ser separado dos demais, e tratamos os resultados de Joshi de
acordo (Seção~\ref{sec:res:h1}). Países com governança estatística comprometida ou
sem padrão nacional publicamente recuperável foram excluídos dos quinze
pré-especificados; a Oceania entrou apenas com a extensão (Austrália), e idiomas
de Classe~0 de Joshi não estão representados, o que registramos como limitação de
escopo.

\subsection{Ampliação da amostra após a análise pré-especificada}
\label{sec:expansao}

Após a análise pré-especificada de 15 países, e \emph{relatado aqui como extensão
post hoc e não como parte do protocolo travado}, a amostra foi ampliada para 25
países sob procedimentos idênticos, para aumentar o poder dos testes de gradiente e
de mecanismo e para multiplicar os pares de idioma dentro do país disponíveis para
H2. O manuscrito relata a amostra de 25 países ao longo de todo o texto; a
família primária recomputada apenas sobre os 15 pré-especificados está tabulada no
material suplementar, de modo que o leitor possa ver o que o protocolo travado
teria concluído. Como a extensão seguiu a observação dos resultados iniciais,
nenhuma conclusão pré-especificada é revisada apenas com base na extensão, e a
ampliação é tratada como checagem de poder e robustez, não como novo teste
confirmatório. Angola, como Nigéria e Argentina, \emph{não} tem padrão nacional de
PM\textsubscript{2,5}, acrescentando casos sem padrão nos quais um modelo fiel deve
se recusar a declarar um limiar inexistente. As dez adições e seus papéis estão
detalhados no suplemento. Duas variantes extras de pergunta por (país, tarefa)
foram geradas para os pares de H2 na fase de desenho, mas nunca foram coletadas
nem pontuadas; a análise de H2 usa a única versão em idioma nativo de cada prompt
em inglês.

\subsection{Covariáveis de país}
\label{sec:covariaveis}

O Índice de Desenvolvimento Humano (PNUD 2023--24) e o PIB per capita indexam a
posição de desenvolvimento e entram na análise de mecanismo como covariáveis de
ajuste. As exposições de representação no corpus se dividem segundo uma distinção
que a análise de mecanismo torna explícita, porque os dois canais se aplicam a
efeitos diferentes. As medidas de \emph{corpus da língua} quantificam quanto texto
de pré-treinamento existe \emph{em} um idioma e são o mecanismo candidato para a
penalidade de idioma nativo: contagem de tokens do mC4~\citep{xue2021mt5}, tamanho
do corpus OSCAR-2301 (corpus apresentado por \citealp{abadji2022oscar}; os
valores por idioma vêm do dataset card do OSCAR-2301, cujos totais diferem da
versão 22.01 descrita naquele artigo) e a classe de
Joshi~\citep{joshi2020state}. As medidas de \emph{corpus do país} quantificam
quanto o corpus enciclopédico fala \emph{sobre} um país, independentemente do
idioma, e são o mecanismo candidato para a lacuna geográfica: tamanho do artigo da
Wikipédia inglesa, número de edições linguísticas com artigo sobre o país e número
de declarações do Wikidata sobre a entidade do país.

As duas famílias não são intercambiáveis. Como a lacuna geográfica é medida
\emph{em inglês}, com o idioma mantido constante entre países, uma lacuna
Norte/Sul residual não pode ser efeito de corpus da língua. Um ponto merece ser
dito de forma direta, porque os papéis se confundem com facilidade: a Wikipédia e
o Wikidata nunca são gabarito aqui, eles são a variável \emph{explicativa}. Aquilo
contra o que uma resposta é pontuada é sempre um registro oficial
(Seção~\ref{sec:gabarito}); tivéssemos pontuado contra uma enciclopédia, o estudo
mediria concordância com uma enciclopédia, que não é a pergunta que uma
administração pública enfrenta.

\subsection{Seleção de modelos}

Catorze LLMs foram avaliados em cinco camadas de acesso, abrangendo cerca de duas
ordens de magnitude em contagem de parâmetros e quatro geografias de dados de
treinamento (Estados Unidos, China, Europa, Brasil): fronteira de pesos abertos
(5), pesos abertos médios (3) e pequenos ou regionais (2), fechados acessíveis (3)
e fechados de fronteira (1). Apenas três dos modelos de pesos abertos foram
servidos localmente (Ollama, tags fixadas: Phi-4~14B, Qwen3~14B, Cabra-Mistral~7B);
os outros sete foram servidos por endpoints hospedados (camada gratuita da Groq
para cinco, APIs da DeepSeek e da Cohere para dois), e os quatro modelos fechados
pelas APIs de seus fornecedores. A camada pequena inclui um modelo ajustado por instrução em português
brasileiro (Cabra-Mistral~7B~v3), testado contra um modelo aberto de treinamento
global de escala equivalente para H3. Cada modelo foi consultado com tag fixada
idêntica ao longo da janela de coleta. O roster completo está na Tabela~S1 do
suplemento.

Comparamos \emph{modelos como servidos} por uma pilha de acesso fixa, e não pesos
idealizados. Infraestrutura do provedor, configuração de serviço e camada de API
fazem parte do que cada rótulo denota, e uma lacuna residual poderia em princípio
refletir a pilha em vez dos pesos; mantemos isso constante fixando tags de versão e
uma configuração de amostragem fixa por provedor ao longo de uma janela fixa
(Seção~\ref{sec:coleta}).

\subsection{Domínio e taxonomia de tarefas}
\label{sec:tarefas}

O domínio único é a política de poluição do ar, escolhido porque se ancora em
gabarito oficialmente publicado e verificável e se mapeia a atividades concretas de
decisão de um gestor ambiental público. As cinco tarefas não foram escritas como um
questionário genérico de conhecimento: cada uma corresponde a uma pergunta que um
gestor tem de responder no curso do trabalho ordinário, e cada uma foi redigida de
modo que uma resposta errada carregasse consequência administrativa, e não apenas a
perda de um ponto. A Tabela~\ref{tab:tarefas} as traz com o prompt tal como foi
emitido aos modelos, usando o Brasil como instância trabalhada, ao lado da decisão
que cada uma representa.

\begin{table}[htbp]
\centering
\caption{As cinco tarefas, o prompt como emitido e a decisão que cada uma
representa. Os prompts são mostrados para a instância Brasil sob a persona
neutra; cada país recebe os mesmos cinco modelos com suas próprias entidades
substituídas (cidade, órgão, regulamento). Os prompts foram emitidos em inglês e
são aqui traduzidos para leitura; o texto original consta do conjunto liberado.}
\label{tab:tarefas}
\footnotesize
\setlength{\tabcolsep}{5pt}
\renewcommand{\arraystretch}{1.0}
\begin{tabular}{@{}>{\raggedright\arraybackslash}p{0.11\linewidth}>{\raggedright\arraybackslash}p{0.54\linewidth}>{\raggedright\arraybackslash}p{0.28\linewidth}@{}}
\toprule
Tarefa & Prompt emitido, em tradução (Brasil) & Decisão que representa \\
\midrule
T1 \newline Padrão técnico
& ``Qual é o padrão nacional vigente de média anual de qualidade do ar para
PM\textsubscript{2,5} no Brasil? Indique o regulamento que o estabelece e se é
meta intermediária ou final. Limite-se a 1--3 frases.''
& Citar o limite vinculante em condicionante de licença, auto de infração ou
parecer técnico. O valor errado vicia o ato. \\
\addlinespace
T2 \newline Dado factual local
& ``Segundo a agência ambiental nacional ou subnacional, qual foi a concentração
média anual de PM\textsubscript{2,5} em São Paulo, Brasil, no ano mais recente
reportado? Forneça o valor em $\mu$g/m\textsuperscript{3}, a estação ou rede de
referência e a fonte.''
& Estabelecer a linha de base contra a qual um plano local de qualidade do ar ou
um relatório público é escrito. \\
\addlinespace
T3 \newline Síntese de evidência em saúde
& ``Resuma a evidência epidemiológica sobre mortalidade atribuível à exposição
ambiental a PM\textsubscript{2,5} no Brasil na última década. Inclua estimativa
de óbitos anuais, causas principais e ao menos uma fonte autoritativa.''
& Justificar por que a qualidade do ar deve ser priorizada frente a demandas
concorrentes pelo mesmo orçamento. \\
\addlinespace
T4 \newline Instrumentos de política
& ``Liste os principais programas nacionais ou instrumentos legais de
monitoramento e controle da poluição do ar ambiente no Brasil, com o órgão
executor de cada um.''
& Identificar de qual programa lançar mão e qual órgão detém a competência. \\
\addlinespace
T5 \newline Recomendação aplicada
& ``Uma região metropolitana do Brasil enfrenta episódio previsto com
PM\textsubscript{2,5} projetado para exceder 75~$\mu$g/m\textsuperscript{3} em
24 horas. Recomende três ações de política de curto prazo (implantáveis em 48
horas) para reduzir a exposição populacional, cada uma justificada por evidência
de saúde pública ou viabilidade operacional.''
& Escolher a resposta operacional a um episódio agudo, sob pressão de tempo. \\
\bottomrule
\end{tabular}
\end{table}

As tarefas estão ordenadas por quanto da resposta é fixado por fonte oficial. T1,
T2 e T4 admitem gabarito verificável e carregam o sinal primário de acurácia; T3 e
T5 são tarefas de geração aberta pontuadas por rubrica, em que respostas
competentes legitimamente diferem. É essa ordenação que torna a dimensão de tarefa
informativa em vez de decorativa: ela permite perguntar se os modelos falham de
maneira uniforme ou falham especificamente onde é preciso recuperar um valor
nacional publicado, que é a pergunta sobre a qual os resultados se apoiam. O
mapeamento de país por tarefa das classes de fonte oficial e do estatuto de
verificabilidade está no registro de gabarito (Seção~\ref{sec:gabarito}), e o
conjunto completo de prompts dos 25 países é liberado com os dados.

\subsection{Manipulação de persona}
\label{sec:persona}

A persona é um fator dentro do prompt com dois níveis, e a manipulação não é
arbitrária: atribuir um papel ao modelo é a prática que os próprios fornecedores
documentam e recomendam (Seção~\ref{sec:related}). A condição neutra apresenta a
consulta como pedido simples de informação; a condição public\_manager\_env
antepõe um enquadramento de papel fixo identificando quem pergunta como autoridade
municipal ou regional de gestão ambiental, mantendo a pergunta substantiva
idêntica. Para a T1 do Brasil, o enquadramento é a frase única \emph{``Sou o gestor
ambiental municipal de São Paulo, Brasil, e preciso informar tomadores de
decisão.''} Um único modelo é aplicado de forma uniforme entre países e tarefas,
traduzido quando o idioma do prompt não é o inglês, de modo que qualquer diferença
entre condições seja atribuível ao enquadramento e não à redação da pergunta. A
ordem das personas dentro de cada célula é aleatorizada, e uma checagem de
manipulação verifica se o enquadramento é reconhecido na resposta. Redação exata,
procedimento de tradução e codificação da checagem estão no suplemento.

Um detalhe de implementação importa para a leitura do resultado. Toda requisição
neste estudo é enviada como turno único de usuário, sem prompt de sistema, de modo
que o enquadramento de papel é anteposto à mensagem do usuário em vez de instalado
como instrução de sistema. Três dos quatro guias o colocam ali; o do Google é
explícito ao dizer que ele cabe \emph{ou} na instrução de sistema \emph{ou} ``no
princípio do prompt do usuário''~\citep{google2026promptingstrategies}, que é
exatamente o canal que testamos. A
escolha é deliberada: a via do prompt de sistema exige acesso por API e um
desenvolvedor para configurá-la, enquanto os servidores públicos cujas decisões
motivam este estudo alcançam esses modelos pela interface de conversa de consumo,
na qual o único canal disponível é a mensagem que digitam. Testamos, portanto, o
enquadramento de papel tal como o setor de fato o executa, e a
Seção~\ref{sec:limites} declara o que permanece não testado. A persona é a base de
H6.

\subsection{Construção e validação dos prompts}

Os prompts foram redigidos para o domínio pela equipe de pesquisa com supervisão de
domínio e validados contra o registro de gabarito antes de entrar no conjunto
confirmatório de estímulos, com a exigência de que toda afirmação factual de uma
entrada do registro resolvesse para uma URL de fonte oficial ou um DOI revisado por
pares antes que o prompt correspondente fosse admitido. O plano de análise previa um
piloto de 50 prompts pontuado por dois avaliadores humanos independentes
($\alpha\geq0.70$) antes de escalar o desenho. \emph{Esse piloto não foi rodado}, e
não existe $\alpha$ de avaliador para este estudo; a confiabilidade da rubrica é
evidenciada, em vez disso, sobre os próprios dados confirmatórios, por um painel de
três fornecedores sobre cada item que exigiu juiz (Seção~\ref{sec:res:confiab}), que
devolve $\alpha=\nalpha$. Relatamos isso como desvio do plano, e não como limiar
atingido.

Todos os prompts foram redigidos em inglês e, para um subconjunto estratificado de
países, adicionalmente em um idioma nativo relevante (Seção~S3 do suplemento). As
versões em idioma nativo foram produzidas por um tradutor LLM (Claude Sonnet~4.6)
instruído a preservar o significado, o enquadramento de persona, nomes próprios e
siglas, com o gabarito factual mantido em inglês, já que o valor-alvo é invariante
ao idioma. Não empregamos tradutores humanos nativos, o que permanece uma limitação
(Seção~\ref{sec:limites}); auditamos, porém, as versões após o fato,
retrotraduzindo todos os 90 prompts nativos com um modelo diferente e tendo dois
outros modelos comparando cada um contra o original em inglês
(Seção~\ref{sec:res:h2}).

\subsection{Gabarito e registro de gabarito}
\label{sec:gabarito}

O gabarito é um conjunto de quatro registros específicos por tarefa, liberados
com os dados, e descrevemos cada um como foi construído, não como foi planejado
(a Tabela~S14 do suplemento registra os desvios). Para T1 a chave é o valor anual
de PM\textsubscript{2,5} em vigor durante a janela de coleta, tomado do próprio
instrumento nacional (lei, decreto, regulamento ou diário oficial) e armazenado
com a fonte, um trecho literal e uma marca de status. Vinte e um países têm chave
numérica. Três (Angola, Argentina, Nigéria) não têm padrão nacional de
PM\textsubscript{2,5}, de modo que a resposta correta é dizê-lo, e um valor
afirmado é pontuado como fabricado. O Egito não tem chave alguma, porque o texto oficial em árabe não pôde ser
recuperado, e fica, portanto, fora de T1 (56 células). Bangladesh é pontuado
contra os 35~$\mu$g/m\textsuperscript{3} de suas Air Pollution (Control) Rules
2022, que substituíram os 15 do diário oficial de 2005; uma versão anterior do
registro trazia o valor de 2005, e a correção está registrada com os arquivos
liberados.
Onde um instrumento fixa uma escada de valores em etapas (as cinco etapas do
CONAMA no Brasil, o limite europeu de 25 com 10 a partir de 2030, o limite de 20
do Reino Unido com meta inglesa de 10 até 2040, os 9,0 dos Estados Unidos que
sucederam os 12,0, os 35 de Bangladesh que sucederam os 15), o registro
armazena a escada inteira além do valor em vigor,
para que citar uma etapa vizinha possa ser separado de fabricação
(Seção~\ref{sec:res:t1}). Os CAAQS do Canadá e o valor filipino são objetivos, e
não limites juridicamente vinculantes; são os únicos valores nacionais de
referência que esses países publicam e são tratados como chave, o que o texto
anota onde importa.

Para T2 a chave é a concentração média anual de PM\textsubscript{2,5} mais recente
que a base WHO Ambient Air Quality Database (v6.1, 2024) registra para a cidade
do prompt, compilando os valores reportados por agências nacionais e
subnacionais; qualquer ano reportado é aceito dentro de $\pm 20\%$, de modo que
um modelo que cite um valor da agência de outro ano não é penalizado pelo ano.
Duas cidades não têm entrada de PM\textsubscript{2,5} (Luanda, Cairo) e essas
células ficam com o juiz. Para T3 a chave é a estimativa do WHO Global Health
Observatory de mortes atribuíveis à poluição do ar ambiente (indicador AIR\_41,
2021), com o intervalo de incerteza alargado em metade de cada lado como faixa
aceita, mais a ordem publicada das causas; o indicador cobre a poluição do ar
ambiente como um todo, enquanto o prompt pergunta por PM\textsubscript{2,5}, e um
modelo que cite outro estudo de carga (o GBD, uma coorte nacional) é pontuado
pelo juiz quanto a se sua cifra e fontes são defensáveis, não contra o número da
OMS. Para T4 o conjunto de referência é o catálogo UNEP Global Air Pollution
Legislation de cada país, com correções manuais onde está desatualizado. O
registro anota fonte, URL, data de acesso e status por entrada, mas não tem campo
de validador humano nem hash por documento: a verificação foi documental, pelos
autores, contra o texto primário. Entradas sem fonte oficial equivalente são
excluídas da comparação adjudicada por código naquela célula, de modo que gabarito
ausente nunca é pontuado como erro do modelo.

\subsection{Coleta de dados}
\label{sec:coleta}

As chamadas foram executadas a temperatura~0.3 com semente fixa onde o provedor a
honra (os três modelos servidos via Ollama), com 2 replicações por combinação de
prompt, modelo e persona, dentro de uma janela fixa para minimizar confundimento
por atualização de modelo; as acurácias relatadas são, portanto, estimativas
pontuais no tempo. Duas exceções são declaradas: a família GPT-5 rejeita qualquer
temperatura diferente da padrão, de modo que GPT-5 e GPT-5-mini rodaram a
temperatura~1.0; e nenhum provedor oferece configuração que torne a saída
determinística, de modo que duas réplicas a 0.3 delimitam, sem remover, a
variância intramodelo (o suplemento compara os efeitos principais com essa
variância). Toda resposta foi
armazenada com metadados completos da requisição --- tag fixada, timestamp,
contagem de tokens, motivo de término, condição de persona --- sob checksums
SHA-256. Gates de qualidade excluíram respostas truncadas e respostas em idioma
divergente do da requisição, ambas retidas e analisadas separadamente.

A cobertura por célula não é uniforme: endpoints gratuitos e com limite de taxa
devolveram menos respostas admissíveis do que as APIs medidas, de modo que o número de
células pontuadas com prompt em inglês varia de 261 a 500 por modelo (777 para
o modelo com mais pares em idioma nativo, quando estes são contados). Relatamos o $N$ realizado de cada
modelo ao lado de sua pontuação (Tabela~\ref{tab:conf-modelo}) em vez de descartar
silenciosamente células sub-cobertas, e a matriz completa de cobertura, com as
razões de cota e exclusão de cada falta, está no suplemento.

\subsection{Medidas}

A pontuação atribui a cada tarefa o instrumento mais confiável que ela admite.

\emph{Onde a resposta é um número em registro oficial, o código decide.} Para
T1, T2 e T3, o valor devolvido é extraído e comparado com o registro por um
pontuador determinístico, com guardas de plausibilidade que rejeitam extrações
implausíveis --- um ano de quatro dígitos lido como contagem de óbitos, uma
população nacional lida como cifra de mortalidade, um valor de 24 horas lido como
anual. O extrator normaliza dígitos devanágari, vírgulas decimais e unidades
escritas por extenso, e seu comportamento nos casos que importam para a
comparação entre idiomas (uma vírgula decimal, um numeral em hindi, um valor
diário entre parênteses ao lado do anual) é fixado por uma suíte de testes
unitários liberada com o código; uma auditoria manual estratificada das respostas
em idioma nativo que ele classificou como sem valor está no suplemento. Comparar
um valor com uma chave é exato, e um modelo de linguagem a quem se pede a mesma
aritmética só consegue reproduzi-la com ruído. O código resolve \ncodecovtone\% das
células de T1 na análise (\ntonecodecellspt{} de \ntonecellspt), \ncodecovttwo\% de T2 e \ncodecovtthree\% de T3; uma
resposta de T1 sem valor extraível conta como não tê-lo devolvido. Para os três
países sem padrão, o código pontua se a resposta afirma que não existe padrão
nacional (correto), afirma um valor (fabricado) ou se recusa a responder (nem um
nem outro), e a comparação primária entre tiers é reportada com e sem eles.

\emph{Onde se exige julgamento, um painel de três fornecedores decide pela média.}
O resíduo de T2 e T3, com toda a T4, é pontuado por Gemini~2.5~Pro, Claude
Sonnet~4.6 e DeepSeek-V3, escolhidos por diversidade de fornecedor para que a
correlação intrafornecedor não infle a concordância aparente. T5, a recomendação
aberta, mantém as pontuações do juiz original da coleta; esse juiz foi o
GPT-5-mini para os quinze países pré-especificados e o Claude Haiku~4.5 para os
dez acrescentados na extensão, e como a extensão contém sete dos dez países do
Norte Global, juiz e tier estão confundidos em qualquer efeito que ainda repouse
sobre essas pontuações. Foi por isso que T1 passou para o código: nas pontuações
originais o déficit do Sul Global em T1 era \njtgsgpt{} contra \njtgngpt{} sob o GPT-5-mini,
mas \njtgshaiku{} contra \njtgnhaiku{} sob o Claude Haiku, uma diferença entre juízes tão grande
quanto a diferença entre tiers (suplemento). T5 não carrega efeito de tier em
nenhum dos dois instrumentos. A confiabilidade é relatada como $\alpha$ de
Krippendorff, o ICC$(2,1)$ de juiz único e, como quantidade operativa, o
ICC$(2,3)$ da média do painel --- um coeficiente de efeitos aleatórios de dois
fatores e concordância absoluta, de modo que as severidades distintas dos juízes
contam contra ele --- computados sobre cada item que o painel pontuou em vez de
sobre um subconjunto de calibração. Um membro do painel (DeepSeek-V3) é também
um modelo avaliado, e testamos o autofavorecimento diretamente. Nenhuma camada
humana de padrão-ouro foi coletada; o gabarito está ancorado, em vez disso, em
fontes oficiais primárias (Seção~\ref{sec:gabarito}).

\emph{A unidade de análise é a célula pontuada, não a linha de log.} A
duplicação existe em dois níveis, e cada um tem uma regra declarada. No nível da
resposta, 21,9\% das chaves (modelo, prompt, réplica) têm mais de uma resposta
armazenada, porque as rodadas de coleta foram retomadas e reexecutadas; a
resposta que entra na análise é a primeira armazenada na ordem dos arquivos de
execução, que é a que o juiz pontuou, e os vereditos de código são computados
sobre essa mesma resposta, para que os dois instrumentos descrevam uma única
resposta. Pontuar todas as respostas armazenadas e tirar a média dos vereditos
muda o veredito de código em \ndupchange\% das células e a média em \ndupmean
(suplemento). No nível do escore, 951 células (11,5\%) carregam mais de um
escore de juiz, porque o juiz foi reexecutado sobre partes do conjunto em datas
diferentes; são medições repetidas do mesmo item, não observações independentes,
e são colapsadas pela média, restando 8.300 células, ou 8.244 uma vez removidas
as 56 células egípcias de T1 sem chave (6.573 com prompt em inglês). Tratar
reexecuções como independentes seria pseudo-replicação: infla o $n$ e dá peso
extra às células que por acaso foram reprocessadas. A escolha é metodológica e
não depende do resultado.

O composto primário combina cinco subcomponentes pré-especificados (Seção~S6 do
suplemento), cada um normalizado para $[0,1]$: acurácia factual contra a entrada do
registro (peso 0.30), completude contextual (0.25), qualidade de citação (0.15),
calibração (0.15) e ausência de alucinação (0.15). Cada efeito principal é
recomputado sob pesos iguais e sob uma ponderação derivada por ACP, e uma conclusão
só conta como robusta se sua direção for estável nas três. O composto é computado
por célula (país, modelo, persona) para viabilizar o teste de diferença em
diferenças de H6.

\subsection{Estratégia analítica}

Os três testes primários são pré-especificados como uma única família (F1) e
corrigidos em conjunto por Bonferroni-Holm. Os testes secundários (F2) são
relatados sem ajuste dentro da família, e os declarados exploratórios (F3) usam FDR
de Benjamini-Hochberg a $q=0.10$~\citep{benjamini1995fdr}, cuja garantia pressupõe
estatísticas de teste independentes.

\emph{H1 (gradiente geográfico, primário)} regride a acurácia composta média por
país sobre o gradiente Joshi/IDH via $\rho$ de Spearman, pré-especificado
unilateral em $\rho\geq0.55$, com um teste de tendência de Mann-Kendall como
complemento obrigatório de não monotonicidade. \emph{H4 (mecanismo de corpus,
primário)} relaciona representação no corpus à acurácia por $\rho$ de Spearman
parcial ajustando para IDH e log do PIB per capita, com contagens da Wikipédia como
proxy pré-especificado. \emph{H6 (persona, primário)} é uma diferença em diferenças
no nível do país, sustentada se a condição de persona reduzir a lacuna do Sul
Global em ao menos 0.05 no composto, com inferência por teste de permutação de
5.000 iterações. \emph{H3 (secundária)} é um contraste unilateral de Mann-Whitney
($\delta$ de Cliff). \emph{H2 e H5 (exploratórias)} são relatadas descritivamente
sob correção F3.

Três análises transversais corroboram a lacuna de camada: um modelo linear misto
gaussiano (MixedLM do \texttt{statsmodels}, REML) com interceptos aleatórios
cruzados de país e modelo, ajustando para IDH centrado; uma reestimação hierárquica
bayesiana (\texttt{pymc}~\citep{pymc}); e E-values para robustez a confundimento não
medido, pela aproximação $\mathrm{RR}\approx\exp(0.91\,d)$ que
\citet{vanderweele2017evalue} prescrevem para efeitos padronizados e cujos
pressupostos ela herda. O mecanismo de corpus é sondado com um modelo
exploratório de mediação (\texttt{semopy}~\citep{semopy}), e a manipulação de
persona com uma checagem de reconhecimento do enquadramento. As análises de
robustez pré-especificadas aparecem na Seção~\ref{sec:res:robustez}.

\subsection{Protocolo do estudo e ética}

O conjunto analítico de dados, após os gates de qualidade mas antes do ajuste dos
modelos, é congelado e hasheado, de modo que a análise confirmatória possa ser
auditada contra a especificação travada em vez de aceita por confiança. Os termos de
liberação estão declarados ao fim do artigo.

Nenhum dado pessoal de usuários finais de LLMs foi coletado; o estudo analisa saídas
de modelo para perguntas de política de relevância pública. Os termos de serviço dos
provedores foram revisados e respeitados, e as respostas são redistribuídas
exclusivamente para pesquisa sob provisões de uso justo e reprodutibilidade,
excluídas de qualquer aplicação comercial.

\section{Resultados}
\label{sec:resultados}

\begin{center}
\fbox{\parbox{0.88\linewidth}{%
\small
Nota sobre a amostra. Os resultados abaixo são relatados sobre a amostra de
\emph{25 países}. O benchmark pré-especificado cobria 15 países e foi estendido
post hoc por dez. Catorze LLMs são pontuados em cinco tarefas, duas personas e
quatro idiomas (inglês mais um idioma nativo do país --- espanhol, português ou
hindi --- para os nove países aplicáveis, o que dá os pares inglês/nativo dentro
do país que testam H2). A família primária recomputada apenas sobre os 15 países
pré-especificados está tabulada no material suplementar, e nenhuma conclusão
pré-especificada é revisada apenas com base na extensão; onde as duas amostras
divergem em uma hipótese de que o artigo depende, dizemos isso no texto.
A pontuação é adjudicada por código sempre que a resposta é um número conferível
contra registro oficial, e pela média de um painel de juízes de três fornecedores
nos demais casos (Seção~\ref{sec:res:confiab}); o gabarito de T1 está ancorado
no instrumento nacional em vigor para 24 países e não tem chave para o Egito,
que fica fora de T1 (Seção~\ref{sec:res:proveniencia}).
}}
\end{center}

\subsection{Amostra e pontuação}
\label{sec:res:amostra}

O corpus analisado compreende \ncellspt{} células pontuadas no composto $[0,1]$, das
quais \ncellsenpt{} são respostas a prompts em inglês nos 25 países; o restante são os
pares em idioma nativo dentro do país usados para H2.

A pontuação usa o instrumento mais confiável que cada tarefa admite. Para T1, T2
e T3, em que a resposta correta é um número em registro oficial, o veredito é
computado por código contra esse registro. O código resolve \ncodecovtone\% das células
de T1, \ncodecovttwo\% de T2 e \ncodecovtthree\% de T3; o resíduo, com toda a T4, é pontuado pela
média de um painel de três fornecedores. T5 mantém as pontuações do juiz da
coleta. No total, \ncodecellspt{} células carregam pontuação de código ou de painel.

A média do painel substitui o desenho de juiz único do protocolo original, no qual
o juiz da coleta (GPT-5-mini para os quinze países pré-especificados, Claude
Haiku~4.5 para os dez da extensão) estava ele próprio entre os modelos avaliados
e, como a extensão contém sete dos dez países do Norte Global, estava confundido
com o tier. Duas medições
justificam a mudança: os juízes diferem acentuadamente em severidade (composto
médio de \njudgegemini{} para o Gemini, \njudgeclaude{} para o Claude e \njudgedeepseek{} para o DeepSeek), de modo
que um juiz único fixa um nível arbitrário, e a confiabilidade de juiz único é de
apenas ICC$(2,1)=0.56$ contra ICC$(2,3)=0.79$ da média do painel. Um membro do
painel (DeepSeek-V3) é também um modelo avaliado, então checamos o
autofavorecimento: ele pontua as próprias saídas de forma \emph{mais} severa em
relação aos outros dois juízes ($-0.177$) do que pontua as de qualquer outro modelo
($-0.123$ a $-0.160$).

\subsection{Acurácia por modelo}
\label{sec:res:modelo}

A Tabela~\ref{tab:conf-modelo} ordena os 14 modelos pelo composto médio. Sistemas
de fronteira lideram e modelos pequenos de pesos abertos ficam atrás. Em H3, o
modelo regional em português brasileiro Cabra-Mistral~7B é o \emph{mais fraco} de
todos (\ncabra{} contra \nrestmean{} para os demais;
Mann-Whitney $p\approx 0$, $\delta$ de Cliff $=\nregdelta$). O $\delta$ de Cliff tem
leitura direta: é a probabilidade de uma resposta sorteada de um grupo superar uma
sorteada do outro, reescalada para variar de $-1$ a $+1$; um valor de $\nregdelta$
significa que o modelo regional perde cerca de três dessas comparações para cada
uma que ganha. Essa comparação é contra o roster inteiro, que inclui sistemas de
fronteira; o teste pré-especificado era mais estreito, Cabra contra seu par de
escala equivalente Llama~3.1~8B em prompts em português, e aponta na mesma
direção com margem menor: no Brasil o modelo regional pontua \nhthreebra{} contra \nhthreebrallama{}
($\delta=\nhthreebradelta$, $n=20$ cada), e nos três países lusófonos \nhthreelus{} contra \nhthreeluslama{}
($\delta=\nhthreelusdelta$, $n=60$ cada). O ajuste regional por instrução no patamar de 7B
não recupera o que escala e treinamento de fronteira dão, mesmo no idioma e no
país para os quais foi ajustado, embora com 20 células brasileiras por modelo o
teste estreito seja descritivo.

\begin{table}[h]
\centering
\caption{Acurácia composta por modelo nos 25 países (adjudicada por código onde a
resposta é valor de registro, média do painel nos demais casos; $[0,1]$). $N$ =
respostas pontuadas a prompts em inglês.}
\label{tab:conf-modelo}
\small
% AUTO-GERADO por code/analysis/make_body_tables.py a partir de freeze_all_effects.json - NAO EDITAR
\begin{tabular}{lrr}
\toprule
Modelo & $N$ & Média \\
\midrule
DeepSeek-V3      & 496 & 0.687 \\
GPT-5            & 496 & 0.674 \\
Gemini~2.5~Flash & 486 & 0.634 \\
GPT-5-mini       & 495 & 0.626 \\
Claude Haiku~4.5 & 486 & 0.591 \\
Qwen3 32B        & 494 & 0.562 \\
Command~R+       & 496 & 0.530 \\
Llama~3.3 70B    & 391 & 0.511 \\
Llama~4 Scout    & 496 & 0.498 \\
Qwen3 14B        & 496 & 0.485 \\
GPT-OSS 120B     & 257 & 0.484 \\
Phi-4 14B        & 496 & 0.477 \\
Llama~3.1 8B     & 496 & 0.438 \\
Cabra-Mistral~7B & 492 & 0.351 \\
\bottomrule
\end{tabular}

\end{table}

\subsection{Acurácia por tarefa: o piso de recuperação}
\label{sec:res:tarefa}

A dificuldade depende fortemente da tarefa (Tabela~\ref{tab:conf-tarefa}). A mais
difícil por larga margem é a T1 (recuperação do padrão técnico, \ntasktone), que exige
devolver o padrão nacional de média anual de PM\textsubscript{2,5} em vigor. A
recomendação aberta (T5) tem a maior pontuação (\ntasktfive) porque sua rubrica premia
conselho de política plausível e bem estruturado, que não depende de um fato único.
Agrupando as duas tarefas de recuperação (T1, T2; média \nfloorrecall) contra síntese e
recomendação (T3--T5, \nfloorrest), a diferença é grande (Mann-Whitney $p\approx 0$,
$\delta$ de Cliff $=\nfloordelta$, cerca de duas perdas e meia para cada vitória). Os
LLMs são mais fracos exatamente onde a gestão ambiental aplicada mais precisa
deles: recuperar o limiar regulatório em vigor.

As duas correções da pontuação são elas próprias informativas. A T2 era pontuada
contra uma chave de preenchimento e sobe de 0.368 para 0.473 assim que cada
concentração devolvida passa a ser comparada por código com a faixa autorizada: a
correção moveu a tarefa cuja chave era defeituosa. A T1 sempre teve chave oficial,
mas era pontuada por dois juízes de coleta de severidades distintas, e passá-la
ao código a eleva de 0.367 para \ntasktone{} mantendo-a como a tarefa mais difícil; o
veredito de código também mostra o que os juízes não podiam: mais de um terço
das respostas erradas de T1 (\nladderhit{} de \nincorrectn) cita um valor que está na escada oficial do
país --- um valor superado ou uma etapa futura --- e não um valor inventado
(Seção~\ref{sec:res:t1}). O piso é real, mais raso do que se relatou primeiro, e
concentrado na recuperação normativa (T1).

\begin{table}[h]
\centering
\caption{Acurácia composta por tipo de tarefa (25 países; T1 em 24, sem chave
para o Egito). T1, T2 e T3 são adjudicadas por código onde o valor devolvido é
conferível contra o registro oficial; T4 é pontuada por painel; T5 mantém as
pontuações do juiz da coleta.}
\label{tab:conf-tarefa}
\small
% AUTO-GERADO por code/analysis/make_body_tables.py a partir de freeze_all_effects.json - NAO EDITAR
\begin{tabular}{lrr}
\toprule
Tarefa & $N$ & Média \\
\midrule
T5 (recomendação aplicada)     & 1.320 & 0.769 \\
T3 (síntese de evid.\ em saúde) & 1.329 & 0.548 \\
T4 (instrumentos de política)  & 1.310 & 0.520 \\
T2 (dado factual local)        & 1.335 & 0.473 \\
T1 (padrão técnico)            & 1.279 & 0.393 \\
\bottomrule
\end{tabular}

\end{table}

\subsection{H2 --- perguntar no idioma local reduz a acurácia}
\label{sec:res:h2}

Para os nove países com idioma nativo alvo, cada prompt é renderizado em inglês e
no idioma nativo, mantendo o conteúdo fixo, de modo que o contraste é dentro do
país e dentro do modelo. As replicações são mediadas dentro de cada célula (modelo,
prompt) antes do pareamento, já que parear replicação a replicação inflaria o $n$ e
tornaria o teste anticonservador. Nas 839 células casadas, o prompt em idioma
nativo \emph{reduz} a acurácia composta média em \nnativa{} pontos percentuais (Wilcoxon
$p=\nnativap$). Esse teste trata as células como independentes, e elas
estão aninhadas em nove países e catorze modelos, de modo que reportamos também
a inferência que respeita o aninhamento: um modelo misto sobre as diferenças
pareadas com intercepto aleatório por país dá $\nhtwomixedcountry$~pp (EP \nhtwomixedcountryse,
$p=\nhtwomixedcountryp$), com intercepto aleatório por modelo $\nhtwomixedmodel$~pp
($p=\nhtwomixedmodelp$); com o país como unidade ($n=9$) a penalidade é negativa
em \nhtwocneg{} de nove e tem média $\nhtwocmean\pm\nhtwocse$~pp ($t=\nhtwoct$, $p=\nhtwoctp$); com o modelo
como unidade ($n=14$) é negativa em \nhtwomneg{} de catorze (Wilcoxon $p=\nhtwomwp$). A
penalidade acompanha o nível de recurso do idioma, e as
três são significativas (Tabela~\ref{tab:h2}): hindi $\nnativahi$~pp, espanhol
$\nnativaes$~pp, português $\nnativapt$~pp. A Índia, que pontua mais alto entre os 25 em
inglês, é a que mais cai em hindi.

Como H2 carrega a afirmação central do estudo, tentamos fazê-la desaparecer em vez
de confirmá-la, e ela se sustentou sob todo ataque. Não é impulsionada por nenhum
país (leave-one-out: $\nloocmin$ a $\nloocmax$~pp, todos significativos), por nenhum modelo
($\nloommin$ a $\nloommax$~pp), nem por células atípicas (5\% aparados: $\ntrim$~pp,
$p=\ntrimp$), e sobrevive a toda ponderação do composto. Aparece em todas
as tarefas e em ambas as condições de persona.

Dois ataques respondem às objeções que um leitor cético levanta primeiro. O
primeiro é que juízes-modelo poderiam penalizar prosa não inglesa por estilo, e não
por conteúdo; testamos então um desfecho binário que nenhum juiz toca: a resposta
contém um valor que o extrator consegue comparar com o registro? O extrator lê
numerais devanágari, vírgulas decimais e unidades escritas por extenso, de modo
que um número em hindi ou em espanhol não é penalizado pela notação
(Seção~\ref{sec:metodos}). Prompts em inglês devolvem valor verificável em \nverifen\%
das vezes e em idioma nativo em \nverifnat\%, uma diferença pareada de $\nverifdiff$~pp, com
a versão nativa pior em \nverifworse{} dos \nverifdisc{} pares discordantes (teste de sinais
$p=\nverifp$): hindi $\nverifhi$~pp, espanhol $\nverifes$~pp, enquanto o
português, o mais bem provido dos três, não mostra tal perda ($\nverifpt$~pp). As
respostas em português devolvem um número tão prontamente quanto as em inglês; o
número é errado com mais frequência. No próprio veredito de código (nativo menos
inglês, células adjudicadas por código nos dois idiomas) a penalidade do
português é de $\nhtwocodept$~pp ($n=\nhtwocodeptn$, $p=\nhtwocodeptp$) e a do espanhol de $\nhtwocodees$ pontos
($n=\nhtwocodeesn$, não significativa), $\nhtwocodeall$~pp no geral ($n=\nhtwocodealln$, $p=\nhtwocodeallp$).
Perguntar no idioma local não apenas reduz uma nota julgada; onde o idioma é
escasso, torna o modelo menos propenso a devolver um número conferível, e onde o
idioma é bem provido torna o número devolvido menos vezes certo.

O segundo é que os prompts nativos, produzidos por um tradutor LLM, poderiam
simplesmente perguntar outra coisa. Testamos isso em todos os 90 prompts nativos
únicos --- um censo, não uma amostra ---, retrotraduzindo cada um com um modelo
\emph{diferente} e tendo dois outros modelos auditando o resultado contra o
original em inglês. Nos dois itens que poderiam confundir o efeito, a fidelidade é
completa: nenhum prompt alterou o país e nenhum alterou a quantidade solicitada
(90/90 em ambos). Restringir H2 aos prompts verificados como fiéis o deixa
essencialmente inalterado ($\nfaithful$~pp, $p=\nfaithfulp$) contra $\ndivergent$~pp entre os
divergentes, diferença não distinguível de zero (permutação $p=\nfidelityp$). A fidelidade
da tradução não prevê o tamanho da penalidade, que é o que se espera se a tradução
não a estiver causando.

Separando os pares pelo instrumento que pontuou a célula nativa, a penalidade é
de $\nhtwocode$~pp onde o código decidiu ($n=\nhtwocoden$, $p=\nhtwocodep$), $\nhtwopanel$~pp onde o painel
decidiu ($n=\nhtwopaneln$, $p=\nhtwopanelp$) e $\nhtwoorig$~pp, não significativa, onde permanece a
pontuação do juiz da coleta ($n=\nhtwoorign$, essencialmente T5, a recomendação aberta
sem valor de registro). O subconjunto adjudicado por código é o menor dos dois
porque o código só age onde existe número, de modo que condiciona a que o modelo
tenha produzido valor nos \emph{dois} idiomas e seleciona para fora o modo de
falha mais severo; o subconjunto do painel carrega esse modo de falha, e é por
isso que a leniência do juiz não pode ser tida como cancelada entre idiomas e por
isso o desfecho sem juiz acima é o que encerra a questão.

A implicação aplicada é o oposto da intuição. Um gestor que suspeita que um modelo
em inglês serve mal seu país, e responde perguntando no idioma nacional, torna a
resposta \emph{pior} --- e pior pela maior margem exatamente onde o idioma é menos
representado nos corpora de treinamento.

\begin{table}[h]
\centering
\caption{H2: acurácia nativa menos inglesa por idioma nativo (pareada dentro de
país e modelo, replicações promediadas dentro da célula, amostra de 25 países).}
\label{tab:h2}
\small
\begin{tabular}{llrrr}
\toprule
Idioma nativo (classe Joshi) & Países & $\Delta$ (pp) & $p$ & $N$ células \\
\midrule
Espanhol (5)  & MEX, ARG, PER, COL, CHL & $\nnativaes$  & $\nnativaesp$ & \nnpareses \\
Português (4) & BRA, PRT, AGO           & $\nnativapt$  & $\nnativaptp$  & \nnparespt \\
Hindi (4)     & IND                     & $\nnativahi$ & $\nnativahip$  & \nnpareshi \\
\midrule
\multicolumn{2}{l}{Geral}               & $\nnativasigned$  & $\nnativap$ & \nnpares \\
\bottomrule
\end{tabular}
\end{table}

\subsection{H6 --- a persona não estreita a lacuna geográfica}
\label{sec:res:h6}

O fator persona foi plenamente cruzado, de modo que H6 é testado como diferença em
diferenças. O enquadramento de gestor ambiental público deixa a lacuna Norte/Sul
essencialmente inalterada: $\ndidneutral$~pp sob o enquadramento neutro contra $\ndidmanager$~pp
sob o de gestor, uma DiD de $\ndid$~pp, muito abaixo do limiar pré-especificado de
5~pp; um teste de permutação no nível do país (5.000 permutações do rótulo de
camada) dá $p=\npersonap$. H6 não é sustentada: apresentar a consulta como vinda de um
gestor ambiental público não estreita nem amplifica a desvantagem do Sul Global,
coerente com a evidência de que personas de especialista não melhoram de forma
confiável a acurácia factual~\citep{zheng2024helpful,mollick2025personas}. Uma
checagem de manipulação encontra que 50.2\% das respostas na condição de persona
reconhecem explicitamente o enquadramento, de modo que o nulo não é artefato de o
enquadramento ser universalmente ignorado.

\subsection{H1 --- uma lacuna de camada que se sustenta, um gradiente que não}
\label{sec:res:h1}

A H1 foi pré-especificada em duas partes, e elas se separam com nitidez: uma se
sustenta e a outra não.

A lacuna de camada se sustenta. Os dez países do Norte Global têm média de \naccgn{}
contra \naccgs{} dos quinze do Sul Global, uma lacuna de $\ntiergap$ pontos percentuais cujo
intervalo de confiança bootstrap de 95\% entre países é $\ntiergapci$~pp e
\emph{exclui o zero} (permutação do rótulo de camada, país como unidade,
$p=\ntiergapperm$). Ela é estável à remoção de qualquer país isolado
(leave-one-out $[\nloogapmin,\nloogapmax]$~pp) e à ponderação do composto ($\ngapequal$~pp sob pesos
iguais, $\ngapfactual$~pp apenas no subcomponente factual, onde se amplia). Ela não
depende dos juízes: só nos desfechos adjudicados por código (toda a T1 e a parte
resolvida de T2 e T3), em que nenhum modelo pontuou nada, os países do Norte
Global devolvem o valor de registro em \ncodegn{} das vezes e os do Sul Global em
\ncodegs{}, uma lacuna de $\ncodegap$~pp $\ncodegapci$ (permutação $p=\ncodegapp$), e a razão
de chances agregada é \ncodeor{} $\ncodeorci$. O que quer que o painel acrescente nas
tarefas julgadas, a lacuna já está lá antes de qualquer juiz ser consultado. Como sete dos
dez países do Norte Global entraram com a extensão, reportamos também a lacuna
onde a onda não pode explicá-la: só dentro da extensão (sete do Norte contra
Colômbia, Chile e Angola) ela é de $\ntiergapwave$~pp $\ntiergapwaveci$, e nos quinze
pré-especificados (três do Norte contra doze) de $\ntiergappre$~pp $\ntiergappreci$, uma
estimativa pontual do mesmo tamanho numa amostra pequena demais para
significância por permutação ($p=\ntiergapprep$). A lacuna é propriedade dos dados, não
da onda que os coletou.

\emph{Onde a lacuna vive, e por que o composto a subestima.} O composto faz a média
de cinco tarefas, e a lacuna não se distribui igualmente entre elas. Reestimando o
mesmo contraste como modelo misto binomial em cada tarefa separadamente --- desfecho
codificado como devolver o valor de registro, com interceptos aleatórios cruzados de
país e modelo ---, a desvantagem escala com o quanto a tarefa depende de um fato
específico daquele país (Tabela~\ref{tab:camada-por-tarefa}). No padrão nacional
vinculante (T1), pontuado por código contra o instrumento em vigor, a chance de
o Sul Global devolver o valor de registro é \nort{} da chance do Norte (\nortgn\%
contra \nortgs\% em taxa bruta), e nos instrumentos de política (T4) \nortfour;
na síntese de evidência em saúde (T3), em que a resposta é uma estimativa
internacional e não nacional, a razão sobe para \nortthree; e na recomendação aberta
(T5), a única tarefa sem valor de registro a acertar, a lacuna desaparece
inteiramente. A razão de T1 sobrevive às duas objeções que sua composição convida
(Seção~\ref{sec:res:t1}): deixando de lado os três países do Sul Global sem
padrão nacional ela é \nortexcl{} $\nortexclci$, e nos quinze pré-especificados é \nortpre{}
$\nortpreci$. A T4 é a única tarefa da tabela pontuada inteiramente pelo
painel, de modo que sua razão pode carregar a ignorância do próprio juiz sobre
os instrumentos do Sul além da do modelo; as estimativas só de código acima são
as livres disso. A lacuna não é fraca --- ela
é \emph{concentrada}, e o composto a dilui ao mediar tarefas em que é severa com
uma tarefa em que não existe. Na pergunta que um gestor público mais leva a esses
sistemas, qual é o padrão em vigor aqui, um país do Sul Global é servido a um
terço das chances de um do Norte.

\begin{table}[h]
\centering
\caption{Chances do Sul Global contra o Norte Global de resposta factualmente
correta, por tarefa. O desfecho é acurácia factual $\ge 0.5$: o veredito de
código onde a tarefa é adjudicada por código (toda a T1; a parte resolvível de T2
e T3), a nota do painel ou do juiz nos demais casos, de modo que para T5, que não
tem valor de registro, é a nota factual do juiz. Modelo misto binomial,
interceptos aleatórios cruzados para país e modelo.}
\label{tab:camada-por-tarefa}
\small
\begin{tabular}{llrrrl}
\toprule
Tarefa & A resposta é & NG & SG & RC & Intervalo 95\% \\
\midrule
T1 (padrão técnico)          & um valor nacional      & \nortonegn\% & \nortonegs\% & \nortone & $\nortoneci$ \\
T4 (instrumentos)            & instrumentos nacionais & \nortfourgn\% & \nortfourgs\% & \nortfour & $\nortfourci$ \\
T2 (dado medido local)       & um valor nacional      & \norttwogn\% & \norttwogs\% & \norttwo & $\norttwoci$ \\
T3 (síntese de evid.\ saúde) & estimativa internac.   & \nortthreegn\% & \nortthreegs\% & \nortthree & $\nortthreeci$ \\
T5 (recomendação aplicada)   & nenhum valor de reg.   & \nortfivegn\% & \nortfivegs\% & \nortfive & $\nortfiveci$ \\
\bottomrule
\end{tabular}
\end{table}

\emph{O gradiente monotônico não se sustenta.} Em $n=25$, o $\rho$ de Spearman
entre acurácia e IDH é de $+\nrho$ e não alcança a significância convencional
($p=\nrhop$; Mann-Kendall $p=\nrhomkp$), ficando bem abaixo do limiar
pré-especificado $\rho\geq0.55$, que é o critério contra o qual H1 foi fixada.
No $n=15$ pré-especificado ele está essencialmente ausente ($\rho=+\nrhopre$), e
nenhuma reestimação \emph{leave-one-country-out} alcança o limiar de $0.55$
(faixa $[\nloorhomin,\nloorhomax]$); sob pesos iguais é $\nrhoequal$ e sob pesos de ACP $\nrhopca$. Relatamos, portanto, uma diferença de camada que
conseguimos demonstrar e um gradiente de desenvolvimento que não: a desvantagem
está estabelecida como diferença entre grupos, enquanto sua forma ao longo da faixa
de desenvolvimento permanece exploratória. Sob Bonferroni-Holm na família primária
$\{$H1-gradiente, H4, H6$\}$, nenhum teste rejeita seu nulo.

A ordenação dos países retém grande heterogeneidade ortogonal aos eixos (material
suplementar): a Índia (\nindia, Sul Global) lidera todos os países, Indonésia e
África do Sul superam a maior parte do Norte Global, e Canadá e França ficam abaixo. A
desvantagem é uma diferença de camada sobreposta a variação substancial específica
de cada país, e essa heterogeneidade é ela própria parte da razão de o gradiente não
se resolver.

\subsection{H4 --- representação no corpus: cobertura do país, inseparável do desenvolvimento}
\label{sec:res:h4}

O mecanismo pré-especificado era a representação no corpus de treinamento, com
proxy de contagem de artigos e edições da Wikipédia. Esse proxy é nulo: $\rho$ de
Spearman $=\nhfourwiki$ ($p=\nhfourwikip$) sem controle, e continua nulo parcializando IDH e log
do PIB per capita. Um proxy post hoc de cobertura do país alcança $\rho=\nhfoursite$
($p=\nhfoursitep$) e se atenua após ajuste por IDH (parcial $\rho=\nhfourpartial$, $p=\nhfourpartialp$). No
nível do país, nenhum dos canais fica estabelecido, porque com 25 países cobertura
e desenvolvimento são colineares.

\subsubsection*{Mudar a unidade de variação resolve o que o nível do país não
resolve}

O obstáculo é confundimento, não ausência de sinal, e ele tem saída que não exige
mais países. Cada país responde a cinco tarefas que diferem em quanto a resposta
depende de um fato sobre \emph{aquele} país: T1, T2 e T4 dependem; T3 (uma
estimativa internacional da OMS) e T5 (uma recomendação genérica) não. Os países
pontuam \nhfourdep{} no primeiro bloco e \nhfourind{} no segundo, um déficit de especificidade de
\nhfourdeficit. Se o que limita o modelo é quanto o corpus diz sobre um dado país, esse
déficit deveria ser menor onde a cobertura é maior. Estimar a interação entre
cobertura e dependência da tarefa no nível da resposta, com efeito fixo de país,
dá $\beta=\nhfourbeta$ (EP \nhfourse, $p=\nhfourp$, $n=\ncellsenpt$). O efeito fixo
absorve todo efeito \emph{principal} no nível do país --- desenvolvimento, renda,
idioma oficial ---, mas não a interação deles com a dependência da tarefa, e é aí
que mora o rival: um país mais rico pode simplesmente ter mais de sua regulação
on-line, alimentando tanto os dados de treinamento do modelo quanto a facilidade
de construir a chave.

Quatro checagens separam o sinal de cobertura de artefato
(Tabela~\ref{tab:h4-dentro}). O contraste \emph{mais puro} --- T1, inteiramente
nacional, contra T5, inteiramente genérica --- produz o \emph{maior} coeficiente;
retirar T5 ou T4 o deixa intacto; um placebo que inverte os rótulos devolve a
imagem espelhada exata; e, entrando com o proxy de corpus da língua ao lado dele,
a cobertura sobrevive ($\beta=\nhfourcob$, $p=\nhfourcobp$) e a língua não
($\beta=\nhfourlang$, $p=\nhfourlangp$). Como toda resposta aqui foi eliciada em inglês, esse
último resultado era predito e é a parte falsificável do desenho. O rival que ele
não elimina é o desenvolvimento. Cobertura e IDH correlacionam a $\rho=\nrhocobhdi$
entre os 25 países, e quando as duas interações entram no mesmo modelo a ordem se
inverte: IDH$\times$dependência sobrevive ($\beta=\nhfourjhdi$, EP 0.007,
$p=\nhfourjhdip$) e cobertura$\times$dependência não ($\beta=\nhfourjcob$, EP
0.007, $p=\nhfourjcobp$). Dentro dos nove países anglófonos, onde o idioma da pergunta é
o mesmo por construção, a cobertura prediz o déficit na magnitude plena
($\beta=\nhfouranglo$, $p=\nhfouranglop$); isso remove a língua como explicação, não o
desenvolvimento.

A afirmação resultante é, portanto, mais estreita que ``a cobertura causa a
lacuna'', e mais estreita do que a escrevemos primeiro. O que o desenho dentro do
país estabelece é que o déficit específico do país é menor onde o mundo registra
mais sobre o país --- seja esse registro medido como cobertura enciclopédica ou
como desenvolvimento, que nesta amostra são duas leituras da mesma variável. Ele
descarta o canal do corpus da língua; não escolhe cobertura sobre
desenvolvimento, e reportamos o mecanismo como identificado apenas até esse par.

\begin{table}[h]
\centering
\caption{Cobertura do país $\times$ dependência da tarefa, modelo misto no nível da
resposta com efeito fixo de país. $\beta$ positivo = maior cobertura, menor déficit
nas tarefas dependentes do país.}
\label{tab:h4-dentro}
\small
\begin{tabular}{lrl}
\toprule
Classificação das tarefas & $\beta$ & $p$ \\
\midrule
T1, T2, T4 contra T3, T5 (principal) & $\nhfourbeta$ & $\nhfourp$ \\
T1 contra T5 (contraste mais puro)   & $\nhfourtonevsfive$ & $\nhfourtonevsfivep$ \\
Excluindo T5 (teto em 99.7\%)        & $\nhfournotfive$ & $\nhfournotfivep$ \\
Excluindo T4 (mais interpretativa)   & $\nhfournotfour$ & $\nhfournotfourp$ \\
Placebo: rótulos invertidos          & $\nhfourplacebo$ & $\nhfourp$ \\
Com IDH$\times$dependência: cobertura & $\nhfourjcob$ & $\nhfourjcobp$ \\
\quad IDH                            & $\nhfourjhdi$ & $\nhfourjhdip$ \\
\bottomrule
\end{tabular}
\end{table}

A explicação de corpus se mostra com mais nitidez em H2, onde a unidade é o idioma: a
penalidade por idioma cai monotonicamente com o volume de tokens do
mC4~\citep{xue2021mt5} e o tamanho do OSCAR~\citep{abadji2022oscar}, com o hindi, o
menor corpus, carregando a maior penalidade --- ainda que, com três idiomas, isso
permaneça descritivo.

\subsection{H5 --- abertos versus fechados acessíveis}
\label{sec:res:h5}

Modelos fechados acessíveis têm média $\nhfive$~pp acima dos de pesos abertos, contra
o enquadramento otimista de H5. O braço aberto inclui modelos maiores, de modo que
a lacuna não é atribuível apenas à escala, ainda que o DeepSeek-V3 lidere no geral.
H5 é interpretada descritivamente.

\subsection{Proveniência do gabarito, e o que a chave de T1 decide}
\label{sec:res:proveniencia}
\label{sec:res:t1}

Toda chave vem de um registro oficial, ancorada por verificação documental e
reproduzível a partir das URLs citadas: T1 da regulação nacional, T2 da Base de
Dados de Qualidade do Ar Ambiente da OMS v6.1, T3 do indicador AIR\_41 do GHO/OMS,
e T4 do Global Assessment of Air Pollution Legislation do PNUMA. Três países
\emph{não} têm padrão nacional de PM\textsubscript{2,5}, confirmado no texto
oficial --- a Nigéria regula apenas PM\textsubscript{10}, a Argentina não tem valor
federal vinculante, Angola não tem legislação nacional ---, e é aí que um modelo
fiel deve se recusar a declarar um limiar inexistente. O Egito não tem valor recuperável e, portanto, não tem chave de T1. Todos esses
quatro casos são do Sul Global, de modo que a composição da chave é ela própria
uma explicação candidata para a lacuna de T1. Não é: deixando de lado os três
países sem padrão, a chance do Sul Global em T1 é \nortexcl{} $\nortexclci$ da do
Norte, contra \nort{} com eles. A chave também se moveu sob revisão, na direção
que o argumento da Seção~\ref{sec:discussao} prevê: a entrada de Bangladesh
trazia os 15~$\mu$g/m\textsuperscript{3} do diário oficial de 2005 até que um
leitor externo apontasse as Rules de 2022, que fixam 35, e com a chave corrigida
a acurácia de Bangladesh em T1 \emph{cai} de 0.39 para \nbgdt, porque a maioria
dos modelos cita o valor superado. O
registro completo está no suplemento. Wikipédia e Wikidata aparecem apenas como
variável explicativa do teste de mecanismo; nenhuma resposta é pontuada contra
elas.

\emph{A chave em vigor, e a escada ao redor dela.} Vários instrumentos fixam uma
escada de valores em vez de um só: a Resolução CONAMA 506/2024 do Brasil percorre
20, 17, 15, 10, 5~$\mu$g/m\textsuperscript{3} em cinco etapas, das quais 17 está
em vigor desde janeiro de 2025; o limite europeu é 25 com 10 a partir de 2030; o
valor-limite do Reino Unido é 20 com uma meta inglesa de 10 até 2040; os Estados
Unidos passaram de 12,0 para 9,0 em 2024; Bangladesh de 15 para 35 em 2022. A
chave de T1 é o valor em vigor na
janela de coleta, e ela é exigente: nenhum modelo devolve 17 para o Brasil (a
resposta modal é 15, a etapa seguinte), quase nenhum devolve 20 para o Reino
Unido (a maioria devolve a meta de 2040 ou o 25 superado), e menos de um quarto
devolve 9,0 para os Estados Unidos. Das \nincorrectn{} respostas erradas de T1, \nladderhit{} citam um
valor que está em algum degrau da escada oficial do país --- uma etapa ainda não
vigente ou um valor já substituído ---, o que é confusão sobre \emph{qual} valor
vincula, não fabricação, e é a falha que um corte de treinamento produziria. A
distinção não salva a comparação entre tiers, ela a aguça: pontuar qualquer valor
da escada como correto eleva a acurácia do Norte Global em T1 de \nstrictgn{} para \nladdergn{} e
a do Sul Global apenas de \nstrictgs{} para \nladdergs{}, porque os instrumentos escalonados
estão em sua maioria no Norte, e a razão de chances cai para \nortladder{}
$\nortladderci$. Onde o registro oferece vários números plausíveis os modelos os
conhecem; onde oferece um só, com frequência não.

\subsection{Modos qualitativos de falha}
\label{sec:res:qualitativo}

Três modos de falha recorrem por trás do composto, com exemplos literais no
suplemento. Os modelos \emph{fabricam padrões inexistentes} para os três países sem
norma nacional: em \nnsfab{} das \nnscells{} células de T1 pontuadas para Angola, Argentina e
Nigéria a resposta afirma um valor nacional específico, em \nnsabst{} afirma corretamente
que nenhum existe e em \nnsdk{} se recusa a responder, e os valores inventados
discordam entre replicações; \emph{erram o valor vinculante}, respondendo
15~$\mu$g/m\textsuperscript{3} para o Brasil, cuja etapa em vigor é 17 (CONAMA
506/2024), uma confusão de etapa e não uma invenção; e \emph{degradam em idioma
nativo}, como quando o GPT-OSS~120B devolveu
os corretos 40~$\mu$g/m\textsuperscript{3} da Índia em inglês e resposta
\emph{vazia} em hindi.

\subsection{Confiabilidade entre juízes}
\label{sec:res:confiab}

A confiabilidade é relatada sobre a amostra que produz os efeitos, não sobre um
subconjunto de calibração. Todos os 3.190 itens que exigiram juiz foram pontuados
pelos três membros do painel, dando $\alpha$ de Krippendorff $=\nalpha$, ICC$(2,1)$
de juiz único $=\nicc$ e --- a quantidade operativa, já que a análise usa a média
do painel --- ICC$(2,3)=\niccpanel$.

O painel também isola uma regra de projeto que vale enunciar de forma geral. Onde o
prompt de julgamento pedia ao modelo que comparasse uma concentração devolvida
contra uma faixa de tolerância expressa em palavras --- isto é, que fizesse ele
próprio a aritmética ---, a concordância entre juízes nesses itens foi de
$r=-0.04$. Pré-computar a faixa admissível e perguntar apenas se o valor cai dentro
dela elevou a concordância nos mesmos itens para $r=+1.00$, e as 274 pontuações
coletadas sob a redação ambígua foram descartadas e recoletadas. Discordância
aparente entre juízes pode ser artefato do que se pede ao juiz que compute. A
aritmética pertence ao código.

\subsection{Robustez}
\label{sec:res:robustez}

Os dois efeitos sustentados foram testados contra sete ataques e sobrevivem a todos;
o gradiente de desenvolvimento fica abaixo do 0.55 pré-especificado em cada um
deles (seu valor mais alto, \nloorhomaxabs, é uma reestimação leave-one-out), o que é o que
torna seu estatuto exploratório um achado, e não uma omissão. As saídas completas estão no
suplemento.

\emph{Nenhum país isolado, e nenhuma ponderação, carrega o resultado.} A reestimação
\emph{leave-one-country-out} mantém a lacuna de camada em $[\nloogapmin,\nloogapmax]$~pp e a
penalidade de idioma nativo entre $\nloocmin$ e $\nloocmax$~pp. Reponderar o composto não move
as conclusões: sob pesos iguais a lacuna é de $\ngapequal$~pp e a penalidade de idioma de
$\nnativaequal$~pp; restringir à acurácia factual contra o valor de registro as amplia para
$\ngapfactual$ e $\nnativafactual$~pp e aprofunda o piso de recuperação. Os efeitos são mais fortes
onde a medida é mais objetiva. As duas réplicas de cada célula diferem em \nrepdiff~pp
em média (DP intracélula de \nrepsd~pp), de modo que a penalidade de idioma de
$\nnativasigned$~pp, mediada sobre 839 pares, fica a cerca de dez erros-padrão do ruído que
duas réplicas deixam por resolver.

\emph{A lacuna de camada sobrevive a um teste livre de distribuição e a duas
famílias de modelo.} Um teste de permutação embaralha os rótulos Norte/Sul entre os
25 países dez mil vezes e conta com que frequência o acaso sozinho produz uma lacuna
desse tamanho; ele não assume nada sobre a distribuição e trata o país, e não a
resposta, como unidade independente. O acaso supera a lacuna observada em \ntiergappermpct\% das
vezes ($p=\ntiergapperm$). Um modelo hierárquico bayesiano com interceptos aleatórios
cruzados concorda ($\beta=\nbayes$, HDI de 94\% $[\nbayeslo,\nbayeshi]$,
$P(\beta<0)=\nbayesp$), enquanto um modelo misto frequentista devolve a mesma magnitude
e sinal mas o deixa marginal ($\beta=\nmixedbeta$, $p=\nmixedp$), porque cobra do termo de
erro a grande heterogeneidade entre países. Dois dos três testes excluem o zero e os
três concordam em magnitude, de modo que relatamos a lacuna como sustentada e
sensível à especificação, com E-value de \nevalue{} na estimativa pontual e \nevaluelim{} no limite
do intervalo mais próximo do nulo~\citep{vanderweele2017evalue} --- um confundidor
de país de força modesta bastaria para explicá-la, e o leitor tem o direito de pesar
isso.

\emph{A lacuna acompanha a posição no sistema-mundo, não a riqueza nacional.}
Reestimada sob três partições baseadas em desenvolvimento dos mesmos 25 países, a
lacuna é significativa sob a UNCTAD e não significativa sob toda divisão traçada
pelo IDH (suplemento). Isso é coerente com o resto do artigo: o gradiente de IDH é
ele próprio nulo, de modo que uma divisão por IDH não deveria produzir lacuna, e o
fato de a UNCTAD separar os dados onde o IDH não separa aponta para a posição no
sistema-mundo, e não a renda, como variável
operante~\citep{quijano2000coloniality,mohamed2020decolonial}. A leitura é post hoc,
de modo que reivindicamos a lacuna para a distinção Norte/Sul tal como
convencionalmente traçada, e não como afirmação geral sobre desenvolvimento.

\emph{O gradiente nulo é informativo, dentro de limites declarados.} Simulando
amostras de 25 países com correlações conhecidas, o desenho tem \npower\% de poder em
$\rho=0.55$ --- o limiar fixado de antemão --- e efeito mínimo detectável de
$\rho=\nmde$. Um gradiente da magnitude que nos comprometemos a chamar de
substancial está excluído; um fraco não está (poder de \npowerweak\% em $\rho=0.40$). Os dez
países acrescentados post hoc caem \emph{dentro} da faixa de IDH que os quinze
originais já cobriam, de modo que a extensão aumenta a densidade de observações, e
não a amplitude do preditor.

\subsection{Síntese}
\label{sec:res:sintese}

A Tabela~\ref{tab:sintese-hip} mapeia cada hipótese ao seu achado, ao mecanismo
candidato e a um estatuto de evidência explícito, de modo que nenhuma afirmação
seja lida como mais forte do que os dados sustentam. Três efeitos são sustentados e
sobrevivem a todo ataque de robustez da Seção~\ref{sec:res:robustez}; duas
expectativas pré-especificadas não se cumprem; e o mecanismo de corpus é
identificado dentro dos países apenas até o par cobertura-ou-desenvolvimento,
que este desenho não consegue separar.

\begin{table}[!ht]
\centering
\caption{Hipóteses, achados, mecanismos candidatos e estatuto de evidência.
Agrupada pelo que a evidência sustenta, não por número de hipótese.}
\label{tab:sintese-hip}
\footnotesize
\setlength{\tabcolsep}{8pt}
\renewcommand{\arraystretch}{1.15}
\renewcommand{\arraystretch}{1.15}
\begin{tabular}{@{}>{\raggedright\arraybackslash}p{0.235\linewidth}%
                  >{\raggedright\arraybackslash}p{0.315\linewidth}%
                  >{\raggedright\arraybackslash}p{0.235\linewidth}@{}}
\toprule
Hipótese & Achado & Mecanismo \\
\midrule
\multicolumn{3}{@{}l}{\emph{Sustentadas}} \\
\addlinespace[2pt]
H2 idioma           & $\nnativasigned$~pp; hindi $\nnativahi$     & corpus do idioma \\
H3 modelo regional  & $\delta=\nregdelta$, pior de 14   & escala sobre ajuste \\
H1 lacuna de camada & $\ntiergap$~pp; RC \nort{} em T1     & camada \\
\midrule
\multicolumn{3}{@{}l}{\emph{Sustentada, mecanismo não separável}} \\
\addlinespace[2pt]
H4 dentro do país   & $\beta=\nhfourbeta$ sozinha; $\nhfourjcob$ ($p=\nhfourjcobp$) com IDH & cobertura ou desenvolvimento \\
\midrule
\multicolumn{3}{@{}l}{\emph{Não sustentadas}} \\
\addlinespace[2pt]
H1 gradiente        & $\rho=\nrho$, abaixo de $0.55$ & desenvolvimento \\
H4 pré-especificada & proxy de corpus do idioma nulo & corpus do idioma \\
H6 persona          & DiD $\ndid$~pp, $p=\npersonap$      & --- \\
\midrule
\multicolumn{3}{@{}l}{\emph{Apenas descritiva}} \\
\addlinespace[2pt]
H5 aberto/fechado   & $\nhfive$~pp para fechados     & tipo de acesso \\
\bottomrule
\end{tabular}

\vspace{4pt}
\begin{minipage}{0.75\linewidth}
\footnotesize\emph{Nota.} O mecanismo de H2 repousa sobre três idiomas nativos e é
descritivo. A lacuna de camada de H1 é concentrada nas tarefas cuja resposta
depende de um fato nacional. H4 dentro do país descarta o canal do corpus da
língua, mas não separa cobertura de desenvolvimento, e não é identificação causal.
\end{minipage}
\end{table}

\section{Discussão}
\label{sec:discussao}

Propusemo-nos a medir se os LLMs respondem a perguntas de política de poluição
do ar com a mesma confiabilidade para o Sul Global e para o Norte Global, se o
enquadramento do prompt muda a resposta e que mecanismo impulsiona qualquer
lacuna. Em 25 países, 14 modelos, quatro idiomas, cinco tarefas e duas
personas, pontuados contra gabarito oficial por código onde a resposta é um
valor de registro e por um painel de três fornecedores nos demais casos, os
danos que conseguimos demonstrar são práticos: os três remédios a que um órgão
recorreria falham, o mais intuitivo deles torna a resposta materialmente pior,
e os modelos são menos confiáveis na recordação do valor vinculante em vigor.
Uma diferença de tier Norte/Sul persiste, e o mecanismo por trás dela não se
estabelece no nível em que a inferência é sólida.

\subsection{Três remédios intuitivos, todos falhando, e um que prejudica}

O idioma local não ajuda e prejudica materialmente. Perguntar no idioma do
próprio país reduz a acurácia em $\nnativa$~pp, significativamente nos três
idiomas e com o país como unidade de inferência, e a maior penalidade recai
sobre o hindi, o idioma de corpus mais escasso, com a Índia, o país de maior
pontuação entre os 25 em inglês, caindo mais. Essa é a direção que a
hierarquia de recursos prevê~\citep{joshi2020state,xue2021mt5}, e a medição
consegue vê-la porque a unidade é o idioma e não o país; com três idiomas, e
espanhol e português indistinguíveis, a ordenação é descritiva. O efeito não
depende de nenhum país, modelo, tarefa, condição de persona ou instrumento de
pontuação isolado, nem da qualidade da tradução, e se mantém em um desfecho que
não envolve juiz algum: o modelo fica marcadamente menos propenso a produzir um
número conferível, e mais propenso a não devolver nada.

A persona de gestor local não estreita a lacuna (DiD $\ndid$~pp, permutação
$p=\npersonap$, equivalente a zero em até $\pm2$~pp), em consonância com a
evidência de que personas de especialista não melhoram a acurácia
factual~\citep{zheng2024helpful,mollick2025personas}, e seu efeito principal é
uma pequena perda ($\npersonamain$~pp). Esse nulo pesa mais do que pesaria um
nulo sobre um hábito popular, porque o enquadramento de papel é a primeira
coisa que todo grande fornecedor diz aos usuários para fazer; Anthropic, OpenAI,
Google e xAI convergem no mesmo conselho (Seção~\ref{sec:related}), de modo que
uma administração que segue a documentação do fornecedor escreve a frase que
testamos. O que essa documentação reivindica para a frase é controle sobre tom,
formato e enquadramento, e nada nela diz que a técnica que molda o registro não
molda também a confiabilidade factual. Nossa estimativa separa as duas. Uma
checagem de manipulação confirma que metade das respostas assume
explicitamente o papel, de modo que o enquadramento é recebido e muda a voz;
ele não muda o que o modelo sabe. Testamo-lo digitado na mensagem e não
instalado como instrução de sistema, o canal que um servidor público numa
interface de chat tem (Seção~\ref{sec:limites}).

O modelo regional é o pior dos catorze ($\delta=\nregdelta$), e fica atrás do
seu par de escala nos prompts brasileiros em português para os quais foi
ajustado ($\nhthreebra$ contra $\nhthreebrallama$ do Llama~3.1~8B, em 20
células cada; $\nhthreenarrow$~pp nos 30 prompts em português), de modo que
recorrer a um modelo pequeno com ajuste local troca a escala e o treinamento de
fronteira que carregam a acurácia no domínio sem comprar de volta o
conhecimento local que promete. Três limites cercam isso: o teste estreito é
pequeno; só nos valores de registro, reservidos nos mesmos itens em português,
Cabra, seu modelo-base e o Llama~3.1~8B ficam a poucas respostas um do outro
(Seção~\ref{sec:res:robustez}), de modo que a margem do composto se apoia nas
tarefas e nos subcomponentes julgados; e o modelo que pudemos rodar é um ajuste
comunitário de uma base de 7B, não um modelo regional comercial de fronteira do
tipo que um órgão poderia adquirir. Para esses, o achado é uma previsão, e o protocolo liberado é o modo
de testá-la.

Os praticantes implantam os três como correções; nenhum funciona, e o que um
gestor mais provavelmente tenta primeiro causa dano ativo. Para órgãos
ambientais do Sul Global, nenhum atalho barato no nível do prompt ou do modelo
substitui a verificação do valor vinculante contra a fonte oficial.

\subsection{O piso de recordação está sobre os valores que vinculam}

O risco recai de forma desigual. Um gestor que redige prosa sobre política de
qualidade do ar é comparativamente bem servido; um gestor que recupera o ponto
de corte regulatório de que uma decisão depende está na pior célula do desenho
(Seção~\ref{sec:res:tarefa}), e pior ainda se perguntar num idioma de menor
recurso. O piso não é artefato de pontuação: o veredito factual nas duas
tarefas de recordação é decidido por código contra o registro, o contraste se
mantém com todo modelo e todo país como unidade, e um terço das respostas
erradas sobre o padrão nacional cita um valor da própria escada oficial do país
e não um valor inventado. Os modelos conhecem a escada e não sabem qual degrau
vincula.

\subsection{Uma diferença de tier que conseguimos mostrar, e um gradiente que não}

O achado geográfico é uma diferença entre grupos, não uma inclinação. Nos
desfechos que nenhum juiz toca, os países do Sul Global devolvem o valor do
registro $\ncodegap$~pp menos vezes do que os do Norte Global, com o país como
unidade; no composto a lacuna é de $\ntiergap$~pp, do mesmo tamanho dentro da
onda de extensão e nos quinze pré-especificados, e inalterada quando os quatro
membros da UE contam como um. Ela é concentrada: cresce com o quanto a resposta
depende de um fato específico do país e está ausente na única tarefa sem valor
de registro a acertar (Seção~\ref{sec:res:h1}), embora os intervalos por tarefa
com o país como unidade sejam largos e essa decomposição seja exploratória. O
gradiente de desenvolvimento fica abaixo do limiar pré-especificado em toda
leitura, e seu intervalo, $[\nrhocilo,\nrhocihi]$, inclui tanto o zero quanto o
limiar: 25 países nem estabelecem um gradiente nem o descartam. A ordenação dos
países mostra por que os dois resultados diferem. A Índia lidera os 25 países
apesar de ser do Sul Global, Canadá e França ficam baixos, de modo que a
desvantagem é uma diferença de tier sobreposta a idiossincrasia específica por
país que uma correlação de postos sobre 25 países não consegue atravessar.

\subsection{O mecanismo, e por que não se estabelece}

O teste de mecanismo pré-especificado falha, e parte da razão é um defeito no
instrumento. O proxy atribui o tamanho da edição da Wikipédia no idioma oficial
dominante do país, de modo que nove dos nossos 25 países recebem o valor da
Wikipédia inglesa. A variável não mede quanto texto existe no idioma de um país;
mede se o país tem inglês como língua oficial, o que acompanha a história
colonial e não o tamanho do corpus. Relatamos isso porque afeta qualquer estudo
que tenha usado ou venha a usar o mesmo proxy.

Mudar a unidade de variação deu uma associação no nível da resposta na direção
prevista, com as verificações que um efeito real passaria, e ela não sobrevive à
inferência com o país como unidade: agrupado por país, o coeficiente tem
$p=\nhfourclp$, os 25 déficits por país se correlacionam com a cobertura a
$\rho=\nhfourctryrho$ ($p=\nhfourctryp$), e dado o IDH a associação é
$\nhfourctrypartial$. Nesses 25 países, cobertura e desenvolvimento se
correlacionam a $\rho=\nrhocobhdi$, e o desenho não tem alavanca que mova um sem
o outro. Relatamos, portanto, a explicação por cobertura como compatível com os
dados e não estabelecida por eles, e a questão cobertura-ou-desenvolvimento
como aberta.

\subsection{A lição metodológica: a chave importa mais que a amostra}

As mesmas células pontuadas sustentam conclusões substantivas diferentes
conforme o instrumento usado para pontuá-las, e o tamanho dessa diferença é a
lição. Pontuar duas tarefas contra chaves provisórias e delegar a aritmética de
faixa a um juiz produziu um gradiente de desenvolvimento significativo e uma
penalidade de idioma pequena; reconstruir as chaves a partir de registros
oficiais e mover a comparação para código reduziu o gradiente pela metade e
mais que dobrou a penalidade (Seções~\ref{sec:res:h1} e \ref{sec:res:h2}).
Mover o veredito factual da terceira tarefa com valor de registro de dois
juízes de coleta para código mudou sua razão de chances-manchete em um terço,
porque os juízes estavam divididos pela mesma linha dos tiers. A própria chave
continuou se movendo sob escrutínio: a entrada do Reino Unido carregava uma
meta de 2040 no lugar do limite em vigor, e a de Bangladesh carregava o valor
de 2005 depois de as Rules de 2022 o terem elevado, de modo que um modelo que
citava o padrão vigente correto de um país do Sul Global era pontuado como
errado até que um leitor externo o apanhou. O conjunto de dados nunca mudou;
apenas o instrumento. O instinto da área quando um efeito de viés geográfico se
revela frágil é acrescentar países, mas o gabarito merece o escrutínio
primeiro. Uma chave defeituosa não apenas acrescenta ruído; acrescenta ruído
estruturado, porque os defeitos se concentram onde os textos oficiais são mais
difíceis de obter ou foram revisados mais recentemente, que é onde a hipótese
mora. O registro liberado carrega, por isso, um hash de conteúdo (material
suplementar), para que qualquer leitor saiba contra qual versão da chave um
número foi computado. A mesma lição vale para a inferência: com a resposta como
unidade, H4 tinha $p=\nhfourp$; com o país como unidade tem $p=\nhfourctryp$, e
só o segundo descreve o que 25 países conseguem mostrar.

\subsection{Por que o domínio importa}

O material particulado fino está entre os principais fatores de risco
ambiental para mortalidade prematura, e a carga recai de forma desproporcional
sobre o Sul Global: a Organização Mundial da Saúde atribui \num{88548} óbitos
no Brasil em 2021 à poluição do ar ambiente (intervalo de incerteza
\numrange{69477}{108431})~\citep{who2026gho_air41}, e a exposição de curto prazo
continua a produzir mortalidade mensurável nas cidades
brasileiras~\citep{requia2024shortterm}. Quando um gestor delega a um LLM uma
consulta normativa ou uma recomendação para um episódio, uma resposta imprecisa
pode se propagar para um parecer regulatório ou uma resposta de emergência.

Não afirmamos que o modelo é menos confiável onde as consequências são maiores,
porque não medimos isso: a carga de doença não está entre nossas covariáveis de
país, e a única observação no nível do país de que dispomos aponta na direção
oposta, já que a Índia carrega uma das maiores cargas de mortalidade por
PM\textsubscript{2,5} do mundo e é também o país de maior pontuação na nossa
amostra. Se a confiabilidade está desalinhada da necessidade é respondível com
a adição de uma covariável de mortalidade atribuível padronizada por idade, mas
não por este desenho. O arcabouço da colonialidade do
saber~\citep{mohamed2020decolonial} oferece uma interpretação do padrão que
medimos, em que os LLMs são resumos comprimidos de corpora produzidos sob
acesso digital e de publicação desigual e a penalidade de idioma é um elo entre
assimetria de corpus e assimetria de conhecimento; os dados sustentam a
penalidade e a diferença de tier, e a interpretação permanece uma
interpretação.

\subsection{Limitações e oportunidades futuras}
\label{sec:limites}

Cada limite abaixo cerca este desenho, e cada um nomeia o estudo que o
removeria. Nós os listamos na ordem em que os executaríamos.

\emph{Sem braço de recuperação.} Toda resposta vem da memória paramétrica do
modelo. Um gestor que usa um assistente com busca, ou um sistema de recuperação
sobre o diário oficial nacional, enfrenta outro instrumento, e a lacuna e a
penalidade de idioma podem mudar sob ele. Os prompts e registros liberados
permitem rodar um braço de recuperação sobre os mesmos itens, e é o primeiro
estudo que acrescentaríamos.

\emph{O enquadramento de papel foi testado no turno do usuário, não no prompt
de sistema.} O enquadramento de gestor é prefixado à mensagem em vez de
instalado como instrução de sistema, onde três dos quatro guias de fornecedor o
colocam; o da Google admite esse lugar ou o início do prompt do usuário, que é o
que usamos. A escolha segue a população que nos interessa, servidores públicos
na interface de chat, que não têm prompt de sistema a configurar, mas nosso
nulo cerca o canal desse usuário, não o de um desenvolvedor. Se uma persona em
instrução de sistema recupera acurácia factual onde uma digitada não recupera
é um experimento de duas células sobre os mesmos itens.

\emph{Do pontuado por painel ao ancorado em padrão-ouro.} Nenhum especialista
humano repontuou respostas aqui. Três propriedades cercam o que isso custa:
onde a resposta é um valor de registro, o veredito factual é computado por
código; onde é preciso julgamento, a pontuação é a média de três fornecedores
e não de um; e todo efeito relatado é um contraste entre grupos, em que a
severidade sistemática de um juiz se cancela. Restringir à acurácia factual
objetiva fortalece os dois efeitos-manchete, o oposto do que um juiz
permissivo produziria. Dois resíduos permanecem: os quatro subcomponentes não
factuais de uma célula de T1 ainda são os do juiz da coleta, razão pela qual o
veredito de código e não o composto é o desfecho primário de tier, e a
auditoria manual do extrator cobriu apenas T1. Uma validação cega por
especialistas humanos de uma amostra estratificada, cerca de 150 respostas
pontuadas por três avaliadores, que o pipeline liberado exporta diretamente,
converteria um resultado ancorado em fonte oficial num resultado ancorado em
padrão-ouro.

\emph{De um sinal dentro do país a um desenho causal.} Dentro dos países
removemos o confundimento no nível do país sem remover o confundimento
tarefa-por-país, e o sinal não alcança significância no nível do país. Um
desenho que varie a cobertura do país mantendo o desenvolvimento fixo, ou um
estudo longitudinal que acompanhe a acurácia à medida que a pegada
enciclopédica de um país cresce, trataria do primeiro problema; mais países
tratariam do segundo.

\emph{De um gradiente nulo a um teste adequadamente potente.} A diferença de
tier está estabelecida enquanto o gradiente não está, com $n=25$, e o intervalo
do gradiente inclui o limiar. Uma replicação pré-registrada com $n$ maior, que
o conjunto de dados aberto torna barata, decidiria se o gradiente está ausente
ou subpotente.

\emph{De três idiomas nativos à matriz multilíngue completa.} O resultado mais
forte se apoia em três idiomas nativos, de modo que seu mecanismo é relatado
descritivamente ainda que o efeito não seja; para um dos três, o hindi, o
próprio instrumento oficial é publicado em inglês, de modo que a penalidade
ali é uma penalidade pelo idioma do usuário e não pelo idioma do registro. A
objeção da qualidade da tradução está resolvida empiricamente pelo censo de
retrotradução relatado acima; o que permanece em aberto é a amplitude, e
escalar para mais idiomas é a extensão de maior valor que este estudo aponta.

\emph{De um domínio a um programa multidomínio.} Todos os efeitos vêm de um
domínio aplicado de alto risco e de uma família de corpus (Wikimedia), e
confinamos a linguagem mecanicista de acordo. A mesma suíte de tarefas e o
mesmo método de registro de gabarito se transferem diretamente a outros
domínios regulados e ancorados em fonte, como qualidade da água, segurança
alimentar e política de vacinação.

\emph{A pilha do modelo servido como unidade de análise.} Avaliamos modelos
tal como servidos por configurações de acesso fixas, que é o que um gestor
público consulta; uma lacuna residual poderia em princípio refletir a pilha em
vez dos pesos. A verificação de serviço do modelo regional
(Seção~\ref{sec:res:robustez}) cobre uma pilha; um desenho controlado com
pesos abertos idênticos em várias pilhas separaria as duas para todas elas, o
que a proveniência por chamada que liberamos permite.

\subsection{O que este estudo contribui}

As três correções a que os praticantes recorrem primeiro falham juntas: o
idioma local prejudica, o enquadramento de persona não faz nada, e o modelo de
7B com ajuste regional é o pior dos catorze. O resultado metodológico é que,
nesta literatura, a chave de gabarito, e o nível em que a inferência é feita,
importam mais que o tamanho da amostra. E o desenho separa o que 25 países
conseguem demonstrar do que não conseguem, relatando o segundo conjunto como
aberto e não como sustentado.

\section{Conclusão}
\label{sec:conclusao}

Modelos de linguagem de grande porte estão se tornando infraestrutura de trabalho
para a política de poluição do ar, domínio em que uma resposta imprecisa sobre um
padrão vinculante, uma concentração medida ou uma resposta operacional carrega
consequências diretas de saúde pública no Sul Global, onde a carga de mortalidade
por PM\textsubscript{2,5} é mais pesada~\citep{gbd2019risk,requia2024shortterm}.
Medimos viés geográfico e modulado por persona nesse domínio com um benchmark de
25 países, 14 modelos e quatro idiomas (inglês mais o idioma nativo de nove
países), pontuado contra gabarito ancorado em
fontes oficiais primárias --- por código sempre que a resposta é um valor de
registro, e pela média de um painel de juízes de três fornecedores nos demais
casos.

As mitigações que um órgão pode tentar não funcionam, e a mais intuitiva sai pela
culatra. Perguntar no idioma do próprio país \emph{reduz} a acurácia em \nnativa{} pontos
percentuais, de forma mais acentuada onde o idioma é menos representado nos
corpora de treinamento; a persona de gestor local não estreita a lacuna; e o
modelo de 7B com ajuste regional é o mais fraco dos catorze sistemas. A acurácia também
despenca na recuperação de um valor vinculante, a tarefa que um gestor menos pode
se dar ao luxo de ter errada. Uma lacuna de camada Norte/Sul persiste e é
altamente estruturada: ao devolver o padrão nacional vinculante, a chance de
resposta correta do Sul Global é \nort{} da chance do Norte, e a lacuna encolhe
conforme a resposta depende menos de um fato específico do país, desaparecendo na
única tarefa sem valor de registro a acertar. O gradiente monotônico de
desenvolvimento não atinge significância, de modo que relatamos uma diferença
entre grupos que conseguimos demonstrar e uma forma ao longo da faixa de
desenvolvimento que não conseguimos. Quanto ao mecanismo, o desenho dentro dos
países descarta o canal do corpus da língua para uma lacuna medida em inglês,
mas não consegue separar a cobertura do país do desenvolvimento, e não
reivindicamos mais do que isso.

Para as agências ambientais do Sul Global, a mensagem prática é cautelar. Saídas
de LLM sobre padrões vinculantes e dados medidos devem ser tratadas como rascunhos
conferidos contra o diário oficial, e nenhum contorno barato no nível do prompt ou
do modelo substitui essa verificação --- muito menos trocar para o idioma local,
que por esta evidência piora a resposta. Interpretamos o padrão pelo arcabouço da
colonialidade do saber~\citep{mohamed2020decolonial}: a penalidade recai sobre
os idiomas sub-atendidos nos corpora de treinamento~\citep{joshi2020state}, e
recai com mais força sobre o hindi, o mais escasso dos três; se esses idiomas são
também falados onde a carga de saúde da poluição do ar é maior não é algo que
este desenho meça, e a única observação no nível de país de que dispomos (Índia)
aponta na direção contrária (Seção~\ref{sec:discussao}).

Um resultado metodológico generaliza para além do domínio. Reconstruir duas chaves
de gabarito a partir de registros oficiais e tirar a comparação de faixa do juiz,
passando-a para o código, mudou as conclusões substantivas sobre as mesmas
respostas. Quando um efeito nesta literatura se mostra frágil, a chave de gabarito
merece a auditoria antes da amostra, porque uma chave defeituosa acrescenta ruído
\emph{estruturado} que se concentra exatamente onde o dado oficial é mais difícil
de obter.

Protocolo, prompts, registros de gabarito com fontes oficiais, código de análise e
dados de resposta anonimizados são liberados abertamente, com pipeline de
reprodução em um comando e um gate de auditoria de métodos que recomputa cada
número relatado a partir do dado bruto. O estudo também oferece um modelo
reutilizável para auditorias aplicadas de equidade em LLMs --- gabarito de fonte
oficial, adjudicação por código sempre que a resposta é um número conferível, e
auditabilidade de ponta a ponta --- nos domínios de alto risco do setor público em
que esses sistemas já estão sendo implantados.


% ============================================================

\section*{Disponibilidade de dados e código}

O protocolo completo, os prompts, o código de análise, os registros de gabarito e
as respostas anonimizadas dos modelos serão liberados na publicação em
\href{https://github.com/Roverlucas/artigo-vies-analise-regional}{github.com/Roverlucas/artigo-vies-analise-regional},
sob licenças MIT (código) e CC-BY-4.0 (dados, prompts,
documentação)~\citep{zenodo2026dataset}. A liberação
inclui os registros por resposta a partir dos quais cada efeito relatado é
calculado, de modo que todo número deste artigo pode ser recomputado do dado
bruto em vez de aceito por confiança. O pipeline de análise roda de ponta a ponta
com \texttt{python code/run\_all.py}.

\section*{Protocolo do estudo}

Hipóteses, marco amostral, definição de desfecho, limiares de decisão e correção
de multiplicidade foram fixados em um plano de análise sob controle de versão
antes do início da coleta confirmatória. Esse plano não foi depositado em
registro público, de modo que o leitor não pode verificar sua data de forma
independente do histórico do repositório, e \emph{não} reivindicamos
pré-registro. O escopo entregue também se afastou substancialmente do desenho
planejado, o que o próprio plano estabelecia como condição para o estatuto
confirmatório.

Relatamos o estudo, portanto, em duas camadas. A primeira declara o que foi
especificado de antemão e o que aconteceu com cada item. A segunda relata o que
os dados sugeriram e o plano não antecipou, rotulado como exploratório do começo
ao fim, com intervalos e sem linguagem confirmatória. Os desvios em relação ao
plano estão tabulados no material suplementar.

\section*{Agradecimentos}

Esta pesquisa não recebeu financiamento específico de agências dos setores
público, comercial ou sem fins lucrativos. Os custos de inferência foram pagos
pelos autores; nenhum fornecedor concedeu créditos, descontos ou apoio em espécie
para este estudo.

\section*{Contribuições dos autores}

Segundo a taxonomia CRediT:
\begin{itemize}\setlength{\itemsep}{0pt}
\item Lucas Rover (UTFPR): Conceituação, Metodologia, Software, Curadoria
  de Dados, Análise Formal, Investigação, Visualização, Redação~--~Rascunho
  Original, Administração do Projeto.
\item Vitor de Melo Dominski (Faculdade Descomplica): Curadoria de Dados,
  Software, Análise Formal.
\item Prof.\ Anibal Tavares de Azevedo (Unicamp): Redação~--~Rascunho
  Original, Redação~--~Revisão e Edição.
\item Dr.\ Eduardo Tadeu Bacalhau (UFPR): Validação, Supervisão,
  Redação~--~Revisão e Edição.
\item Dra.\ Yara de Souza Tadano (UTFPR): Validação, Supervisão,
  Redação~--~Revisão e Edição.
\end{itemize}

\section*{Declaração sobre IA generativa na redação}

É preciso separar dois usos distintos, porque apenas um deles é objeto desta
declaração.

\emph{Como objeto e instrumento do estudo.} Modelos de linguagem são o objeto
desta pesquisa e também servem como um de seus instrumentos de medida: GPT-5-mini,
Gemini~2.5~Pro, Claude Sonnet~4.6 e DeepSeek-V3 pontuaram respostas contra o
registro de gabarito, e o Claude Sonnet~4.6 produziu as versões dos prompts em
idioma nativo. Esses usos fazem parte do método, estão descritos na
Seção~\ref{sec:metodos}, e são a razão de o estudo também relatar confiabilidade
entre juízes e uma auditoria de retrotradução dos prompts traduzidos. Onde a
resposta correta é um valor de registro oficial, o veredito é computado por
código, não por modelo.

\emph{Na preparação do manuscrito.} Durante a preparação deste trabalho, os
autores usaram modelos de linguagem para auxiliar na redação e edição do texto e
na implementação do código de análise. Após usar essas ferramentas, os autores
revisaram e editaram o conteúdo conforme necessário e assumem responsabilidade
integral pelo conteúdo da publicação. Nenhum modelo de linguagem consta como
autor, e nenhum modelo contribuiu com a interpretação dos achados ou com as
decisões sobre o que a evidência sustenta.

\section*{Declaração de interesses}

Os autores declaram não haver conflitos de interesse.

\bibliography{references}

\end{document}
